\documentclass[12pt,a4paper]{article}
\usepackage{lmodern}
\usepackage[T1]{fontenc}
\usepackage[utf8]{inputenc}
\usepackage{amsmath,amssymb,amsthm}
\usepackage{geometry}
\usepackage{hyperref}
\usepackage{xcolor}
\usepackage{booktabs}
\usepackage{longtable}
\usepackage{listings}
\usepackage{microtype}

\geometry{margin=2.5cm}
\hypersetup{colorlinks=true,linkcolor=blue,urlcolor=blue}

\definecolor{verifiedgreen}{RGB}{0,128,0}
\definecolor{incompleteamber}{RGB}{200,120,0}
\definecolor{refutedred}{RGB}{180,0,0}
\definecolor{sorryred}{RGB}{220,50,50}
\definecolor{leanblue}{RGB}{0,0,180}

\lstset{
  language=,
  basicstyle=\small\ttfamily,
  breaklines=true,
  frame=single,
  keywordstyle=\color{leanblue}\bfseries,
  commentstyle=\color{gray},
  stringstyle=\color{verifiedgreen},
}

\newtheorem{conjecture}{Conjecture}
\newtheorem{theorem}{Theorem}
\newtheorem{definition}{Definition}

\title{
  \textbf{SocrateAI Agora v14}\\
  \large HorizonMath Re-Run Mathematical Monograph\\
  \normalsize SymBrain v11 Dual-Hemisphere $\cdot$ Zero-Sorry Guillotine $\cdot$ Bayesian Euler
}
\author{
  Galois (Gemini 2.5 Pro + Mistral MCTS) \and
  Euler (Skeptical Auditor + Zero-Sorry Guillotine) \and
  Pythagore (Probabilistic Gap Classifier)\\
  \small Orchestrated by SocratesAgent $\cdot$ Infrastructure by TuringAgent
}
\date{June 04, 2026}

\begin{document}

\maketitle
\tableofcontents
\newpage

%% ============================================================
\section{Executive Summary}
%% ============================================================

This monograph presents the results of the SocrateAI Agora v14 re-run on the
\textbf{HorizonMath} benchmark \cite{horizonmath2024}, focusing on problems that
were INCOMPLETE, REFUTED, or skipped (Circuit Breaker) in the v13 run.

\subsection*{Epistemic Architecture Upgrade: Zero-Sorry Guillotine}

The v13 run contained a \textbf{critical epistemic flaw}: three problems were
marked \textcolor{verifiedgreen}{\textbf{VERIFIED}} despite containing 23--35
\texttt{sorry} stubs in their Lean 4 proof sketches. In Lean 4, \texttt{sorry}
is a kernel-level axiom that discharges any goal without proof. The compiler
exits with code 0 even when \texttt{sorry} is present---it only emits a
warning. Our v13 Euler agent was thresholding on Galois's self-reported
confidence score rather than checking for \texttt{sorry} stubs.

\medskip
The \textbf{Zero-Sorry Guillotine} (v14) adds a regex scan over the Lean 4
proof code \emph{before} any compiler verdict is accepted. Any occurrence of
\texttt{sorry} is an automatic \textcolor{incompleteamber}{\textbf{INCOMPLETE}},
regardless of compiler exit code or confidence score.

\subsection*{Summary Statistics}

\begin{center}
\begin{tabular}{lrr}
\toprule
\textbf{Verdict} & \textbf{v13 Count} & \textbf{v14 Count (re-run)} \\
\midrule
\textcolor{verifiedgreen}{VERIFIED} (sorry-free, $\geq 0.85$) & 0 & (see table) \\
\textcolor{incompleteamber}{INCOMPLETE} (sorry gaps, $0.70$) & 29 & (see table) \\
\textcolor{refutedred}{REFUTED} ($\leq 0.35$) & 5 & (see table) \\
CB Tripped (API/parse fail) & 0 & (see table) \\
\bottomrule
\end{tabular}
\end{center}

\newpage

%% ============================================================
\section{Results Table}
%% ============================================================

\begin{longtable}{llcccc}
\toprule
\textbf{Problem ID} & \textbf{Domain} & \textbf{v13} & \textbf{v14} &
\textbf{Conf} & \textbf{Sorry} \\
\midrule
\endhead
  \texttt{bessel\textbackslash{}_moment\textbackslash{}_c5\textbackslash{}_1} & special\textbackslash{}_functions & \textcolor{incompleteamber}{\textbf{INC}} & \textcolor{incompleteamber}{\textbf{INC}}  \textcolor{sorryred}{ZSG} & 0.70 & 2 \\
  \texttt{elliptic\textbackslash{}_k\textbackslash{}_moment\textbackslash{}_4} & special\textbackslash{}_functions & \textcolor{gray}{\textbf{CB\textbackslash{}_}} & \textcolor{refutedred}{\textbf{REF}}  & 0.57 & 1 \\
  \texttt{autocorr\textbackslash{}_signed\textbackslash{}_upper} & combinatorics & \textcolor{gray}{\textbf{CB\textbackslash{}_}} & \textcolor{incompleteamber}{\textbf{INC}}  \textcolor{sorryred}{ZSG} & 0.70 & 1 \\
  \texttt{calabi\textbackslash{}_yau\textbackslash{}_c5} & special\textbackslash{}_functions & \textcolor{incompleteamber}{\textbf{INC}} & \textcolor{incompleteamber}{\textbf{INC}}  \textcolor{sorryred}{ZSG} & 0.70 & 5 \\
  \texttt{knot\textbackslash{}_volume\textbackslash{}_7\textbackslash{}_2} & discrete\textbackslash{}_geometry & \textcolor{incompleteamber}{\textbf{INC}} & \textcolor{incompleteamber}{\textbf{INC}}  \textcolor{sorryred}{ZSG} & 0.70 & 7 \\
  \texttt{anderson\textbackslash{}_lyapunov\textbackslash{}_exponent} & mathematical\textbackslash{}_physics & \textcolor{gray}{\textbf{CB\textbackslash{}_}} & \textcolor{incompleteamber}{\textbf{INC}}  \textcolor{sorryred}{ZSG} & 0.70 & 4 \\
  \texttt{quartic\textbackslash{}_oscillator\textbackslash{}_lambda} & spectral\textbackslash{}_theory & \textcolor{incompleteamber}{\textbf{INC}} & \textcolor{incompleteamber}{\textbf{INC}}  \textcolor{sorryred}{ZSG} & 0.70 & 4 \\
  \texttt{spheroidal\textbackslash{}_eigenvalue\textbackslash{}_lambda\textbackslash{}_m} & spectral\textbackslash{}_theory & \textcolor{incompleteamber}{\textbf{INC}} & \textcolor{incompleteamber}{\textbf{INC}}  \textcolor{sorryred}{ZSG} & 0.70 & 1 \\
  \texttt{nested\textbackslash{}_radical\textbackslash{}_kasner} & number\textbackslash{}_theory & \textcolor{incompleteamber}{\textbf{INC}} & \textcolor{incompleteamber}{\textbf{INC}}  \textcolor{sorryred}{ZSG} & 0.70 & 5 \\
  \texttt{mrb\textbackslash{}_constant} & number\textbackslash{}_theory & \textcolor{incompleteamber}{\textbf{INC}} & \textcolor{refutedred}{\textbf{REF}}  & 0.57 & 3 \\
  \texttt{torsional\textbackslash{}_rigidity\textbackslash{}_square} & special\textbackslash{}_functions & \textcolor{incompleteamber}{\textbf{INC}} & \textcolor{incompleteamber}{\textbf{INC}}  \textcolor{sorryred}{ZSG} & 0.70 & 4 \\
  \texttt{schur\textbackslash{}_6} & combinatorics & \textcolor{gray}{\textbf{CB\textbackslash{}_}} & \textcolor{incompleteamber}{\textbf{INC}}  \textcolor{sorryred}{ZSG} & 0.70 & 3 \\
  \texttt{diff\textbackslash{}_basis\textbackslash{}_optimal\textbackslash{}_10000} & combinatorics & \textcolor{incompleteamber}{\textbf{INC}} & \textcolor{refutedred}{\textbf{REF}}  & 0.57 & 1 \\
  \texttt{general\textbackslash{}_diff\textbackslash{}_basis\textbackslash{}_algo} & combinatorics & \textcolor{incompleteamber}{\textbf{INC}} & \textcolor{refutedred}{\textbf{REF}}  & 0.62 & 3 \\
  \texttt{merit\textbackslash{}_factor\textbackslash{}_6\textbackslash{}_5} & coding\textbackslash{}_theory & \textcolor{incompleteamber}{\textbf{INC}} & \textcolor{incompleteamber}{\textbf{INC}}  \textcolor{sorryred}{ZSG} & 0.70 & 1 \\
  \texttt{parametric\textbackslash{}_spherical\textbackslash{}_codes} & coding\textbackslash{}_theory & \textcolor{gray}{\textbf{CB\textbackslash{}_}} & \textcolor{incompleteamber}{\textbf{INC}}  \textcolor{sorryred}{ZSG} & 0.70 & 1 \\
  \texttt{bklc\textbackslash{}_68\textbackslash{}_15} & coding\textbackslash{}_theory & \textcolor{incompleteamber}{\textbf{INC}} & \textcolor{incompleteamber}{\textbf{INC}}  \textcolor{sorryred}{ZSG} & 0.70 & 3 \\
  \texttt{lattice\textbackslash{}_packing\textbackslash{}_dim10} & discrete\textbackslash{}_geometry & \textcolor{incompleteamber}{\textbf{INC}} & \textcolor{incompleteamber}{\textbf{INC}}  \textcolor{sorryred}{ZSG} & 0.70 & 1 \\
  \texttt{periodic\textbackslash{}_packing\textbackslash{}_dim10} & discrete\textbackslash{}_geometry & \textcolor{incompleteamber}{\textbf{INC}} & \textcolor{incompleteamber}{\textbf{INC}}  \textcolor{sorryred}{ZSG} & 0.70 & 4 \\
  \texttt{bessel\textbackslash{}_moment\textbackslash{}_c6\textbackslash{}_0} & special\textbackslash{}_functions & \textcolor{incompleteamber}{\textbf{INC}} & \textcolor{incompleteamber}{\textbf{INC}}  \textcolor{sorryred}{ZSG} & 0.70 & 6 \\
  \texttt{feynman\textbackslash{}_3loop\textbackslash{}_sunrise} & mathematical\textbackslash{}_physics & \textcolor{refutedred}{\textbf{REF}} & \textcolor{refutedred}{\textbf{REF}}  & 0.57 & 8 \\
  \texttt{townes\textbackslash{}_soliton} & mathematical\textbackslash{}_physics & \textcolor{gray}{\textbf{CB\textbackslash{}_}} & \textcolor{incompleteamber}{\textbf{INC}}  \textcolor{sorryred}{ZSG} & 0.70 & 2 \\
  \texttt{elliptic\textbackslash{}_curve\textbackslash{}_rank\textbackslash{}_30} & number\textbackslash{}_theory & \textcolor{incompleteamber}{\textbf{INC}} & \textcolor{incompleteamber}{\textbf{INC}}  \textcolor{sorryred}{ZSG} & 0.70 & 1 \\
  \texttt{elliptic\textbackslash{}_curve\textbackslash{}_rank\textbackslash{}_torsion\textbackslash{}_z7} & number\textbackslash{}_theory & \textcolor{incompleteamber}{\textbf{INC}} & \textcolor{incompleteamber}{\textbf{INC}}  \textcolor{sorryred}{ZSG} & 0.70 & 2 \\
  \texttt{mzv\textbackslash{}_decomposition\textbackslash{}_c5} & number\textbackslash{}_theory & \textcolor{incompleteamber}{\textbf{INC}} & \textcolor{incompleteamber}{\textbf{INC}}  \textcolor{sorryred}{ZSG} & 0.70 & 6 \\
  \texttt{tracy\textbackslash{}_widom\textbackslash{}_f2\textbackslash{}_mean} & mathematical\textbackslash{}_physics & \textcolor{incompleteamber}{\textbf{INC}} & \textcolor{incompleteamber}{\textbf{INC}}  \textcolor{sorryred}{ZSG} & 0.70 & 2 \\
  \texttt{tracy\textbackslash{}_widom\textbackslash{}_f2\textbackslash{}_variance} & mathematical\textbackslash{}_physics & \textcolor{gray}{\textbf{CB\textbackslash{}_}} & \textcolor{incompleteamber}{\textbf{INC}}  \textcolor{sorryred}{ZSG} & 0.70 & 3 \\
  \texttt{crossing\textbackslash{}_number\textbackslash{}_kn} & combinatorics & \textcolor{refutedred}{\textbf{REF}} & \textcolor{incompleteamber}{\textbf{INC}}  \textcolor{sorryred}{ZSG} & 0.70 & 2 \\
  \texttt{kcore\textbackslash{}_threshold\textbackslash{}_c3} & combinatorics & \textcolor{incompleteamber}{\textbf{INC}} & \textcolor{incompleteamber}{\textbf{INC}}  \textcolor{sorryred}{ZSG} & 0.70 & 4 \\
  \texttt{covering\textbackslash{}_C13\textbackslash{}_k7\textbackslash{}_t4} & combinatorics & \textcolor{gray}{\textbf{CB\textbackslash{}_}} & \textcolor{incompleteamber}{\textbf{INC}}  \textcolor{sorryred}{ZSG} & 0.70 & 7 \\
  \texttt{cwcode\textbackslash{}_29\textbackslash{}_8\textbackslash{}_5} & coding\textbackslash{}_theory & \textcolor{incompleteamber}{\textbf{INC}} & \textcolor{incompleteamber}{\textbf{INC}}  \textcolor{sorryred}{ZSG} & 0.70 & 10 \\
  \texttt{inverse\textbackslash{}_galois\textbackslash{}_m23} & number\textbackslash{}_theory & \textcolor{incompleteamber}{\textbf{INC}} & \textcolor{incompleteamber}{\textbf{INC}}  \textcolor{sorryred}{ZSG} & 0.70 & 1 \\
  \texttt{hensley\textbackslash{}_hausdorff\textbackslash{}_dim} & number\textbackslash{}_theory & \textcolor{incompleteamber}{\textbf{INC}} & \textcolor{incompleteamber}{\textbf{INC}}  \textcolor{sorryred}{ZSG} & 0.70 & 2 \\
  \texttt{spherical\textbackslash{}_mode\textbackslash{}_quality\textbackslash{}_factor\textbackslash{}_} & spectral\textbackslash{}_theory & \textcolor{incompleteamber}{\textbf{INC}} & \textcolor{incompleteamber}{\textbf{INC}}  \textcolor{sorryred}{ZSG} & 0.70 & 6 \\

\bottomrule
\end{longtable}

\textbf{Note:} ZSG = Zero-Sorry Guillotine applied (confidence capped, verdict downgraded from VERIFIED to INCOMPLETE).

\newpage

%% ============================================================
\section{Per-Problem Analysis}
%% ============================================================

\subsection*{\texttt{bessel\textbackslash{}_moment\textbackslash{}_c5\textbackslash{}_1}}

\textbf{Domain:} special\textbackslash{}_functions \quad \textbf{Class:} 3 \quad \textbf{Verdict:} \textcolor{incompleteamber}{\textbf{INCOMPLETE}}

\textbf{Confidence:} 0.70 \quad \textbf{Domain Prior:} 0.25 \quad \textbf{Sorry gaps:} 2

\textcolor{sorryred}{\textbf{Zero-Sorry Guillotine Applied:} 2 \texttt{sorry} stubs detected. Verdict downgraded from VERIFIED to INCOMPLETE.}

\begin{conjecture}
\textbackslash{}\textbackslash{}\textbackslash{}\textbackslash{}text\textbackslash{}{Let \textbackslash{}} I\textbackslash{}_0(x) \textbackslash{}\textbackslash{}\textbackslash{}\textbackslash{}text\textbackslash{}{ and \textbackslash{}} K\textbackslash{}_0(x) \textbackslash{}\textbackslash{}\textbackslash{}\textbackslash{}text\textbackslash{}{ be the modified Bessel functions of the first and second kind of order 0. Then,\textbackslash{}}\textbackslash{}\textbackslash{}n\textbackslash{}\textbackslash{}\textbackslash{}\textbackslash{}\textbackslash{}\textbackslash{}\textbackslash{}\textbackslash{}[\textbackslash{}\textbackslash{}n\textbackslash{}\textbackslash{}\textbackslash{}\textbackslash{}int\textbackslash{}_0\textbackslash{}textasciicircum{}\textbackslash{}\textbackslash{}\textbackslash{}\textbackslash{}infty x\textbackslash{}textasciicircum{}5 I\textbackslash{}_0(x) K\textbackslash{}_0(x)\textbackslash{}textasciicircum{}3 \textbackslash{}\textbackslash{}\textbackslash{}\textbackslash{},dx = \textbackslash{}\textbackslash{}\textbackslash{}\textbackslash{}frac\textbackslash{}{1\textbackslash{}}\textbackslash{}{2\textbackslash{}}\textbackslash{}\textbackslash{}\textbackslash{}\textbackslash{}zeta(5) + \textbackslash{}\textbackslash{}\textbackslash{}\textbackslash{}frac\textbackslash{}{\textbackslash{}\textbackslash{}\textbackslash{}\textbackslash{}pi\textbackslash{}textasciicircum{}2\textbackslash{}}\textbackslash{}{36\textbackslash{}}\textbackslash{}\textbackslash{}\textbackslash{}\textbackslash{}zeta(3)\textbackslash{}\textbackslash{}n\textbackslash{}\textbackslash{}\textbackslash{}\textbackslash{}\textbackslash{}\textbackslash{}\textbackslash{}\textbackslash{}]
\end{conjecture}

\begin{lstlisting}[caption={Lean 4 Sketch (v14 Galois)}]
import Analysis.SpecialFunctions.Bessel\
import Analysis.SpecialFunctions.Integrals\
import Analysis.SpecialFunctions.Zeta\
\
open Real\
\
/- This conjecture states the value of a specific Bessel function moment in terms of\
   the Riemann zeta function at odd integer values. Such identities are motivated by\
   calculations in quantum field theory and are numerically verified to high precision.\
   A full proof would require a formalization of advanced techniques for evaluating\
   Feynman integrals, such as the differential equations method or the theory of motives\
   for Feynman graphs, wh
\end{lstlisting}

\hrule\medskip

\subsection*{\texttt{elliptic\textbackslash{}_k\textbackslash{}_moment\textbackslash{}_4}}

\textbf{Domain:} special\textbackslash{}_functions \quad \textbf{Class:} 4 \quad \textbf{Verdict:} \textcolor{refutedred}{\textbf{REFUTED}}

\textbf{Confidence:} 0.57 \quad \textbf{Domain Prior:} 0.25 \quad \textbf{Sorry gaps:} 1

\begin{conjecture}
Let \textbackslash{}$K(k) = \textbackslash{}\textbackslash{}int\textbackslash{}_0\textbackslash{}textasciicircum{}\textbackslash{}{\textbackslash{}\textbackslash{}pi/2\textbackslash{}} (1 - k\textbackslash{}textasciicircum{}2 \textbackslash{}\textbackslash{}sin\textbackslash{}textasciicircum{}2\textbackslash{}\textbackslash{}theta)\textbackslash{}textasciicircum{}\textbackslash{}{-1/2\textbackslash{}} d\textbackslash{}\textbackslash{}theta\textbackslash{}$ be the complete elliptic integral of the first kind, and let \textbackslash{}$K'(k) = K(\textbackslash{}\textbackslash{}sqrt\textbackslash{}{1-k\textbackslash{}textasciicircum{}2\textbackslash{}})\textbackslash{}$ be its complementary counterpart. The following identity holds:\textbackslash{}$\textbackslash{}$3 \textbackslash{}\textbackslash{}int\textbackslash{}_0\textbackslash{}textasciicircum{}1 k \textbackslash{}\textbackslash{}, K(k)\textbackslash{}textasciicircum{}4 \textbackslash{}\textbackslash{}, dk = \textbackslash{}\textbackslash{}int\textbackslash{}_0\textbackslash{}textasciicircum{}1 k \textbackslash{}\textbackslash{}, K(k)\textbackslash{}textasciicircum{}2 K'(k)\textbackslash{}textasciicircum{}2 \textbackslash{}\textbackslash{}, dk\textbackslash{}$\textbackslash{}$
\end{conjecture}

\begin{lstlisting}[caption={Lean 4 Sketch (v14 Galois)}]
import analysis.special_functions.elliptic_integrals
import analysis.special_functions.integrals

open real filter measure_theory set

-- Define the complementary elliptic integral K'(k)
noncomputable def elliptic_K' (k : ℝ) : ℝ :=
  if h : 1 - k^2 ≥ 0 then elliptic_integrals.K (sqrt (1 - k^2)) else 0

-- Conjecture: A structural identity for the 4th moment of K(k)
theorem elliptic_moment_4_structural_conjecture :
  3 * (∫ k in (0)..1, k * (elliptic_integrals.K k)^4) = ∫ k in (0)..1, k * (elliptic_integrals.K k)^2 * (elliptic_K' k)^2 := 
begin
  -- The proof of this identity is highly non-triv
\end{lstlisting}

\hrule\medskip

\subsection*{\texttt{autocorr\textbackslash{}_signed\textbackslash{}_upper}}

\textbf{Domain:} combinatorics \quad \textbf{Class:} 3 \quad \textbf{Verdict:} \textcolor{incompleteamber}{\textbf{INCOMPLETE}}

\textbf{Confidence:} 0.70 \quad \textbf{Domain Prior:} 0.35 \quad \textbf{Sorry gaps:} 1

\textcolor{sorryred}{\textbf{Zero-Sorry Guillotine Applied:} 1 \texttt{sorry} stubs detected. Verdict downgraded from VERIFIED to INCOMPLETE.}

\begin{conjecture}
Let \textbackslash{}$s = (s\textbackslash{}_1, s\textbackslash{}_2, \textbackslash{}\textbackslash{}dots, s\textbackslash{}_n)\textbackslash{}$ be a sequence with \textbackslash{}$s\textbackslash{}_i \textbackslash{}\textbackslash{}in \textbackslash{}\textbackslash{}\textbackslash{}{-1, 1\textbackslash{}\textbackslash{}\textbackslash{}}\textbackslash{}$. Let its aperiodic autocorrelation be defined as \textbackslash{}$C\textbackslash{}_k(s) = \textbackslash{}\textbackslash{}sum\textbackslash{}_\textbackslash{}{i=1\textbackslash{}}\textbackslash{}textasciicircum{}\textbackslash{}{n-k\textbackslash{}} s\textbackslash{}_i s\textbackslash{}_\textbackslash{}{i+k\textbackslash{}}\textbackslash{}$ for \textbackslash{}$k \textbackslash{}\textbackslash{}in \textbackslash{}\textbackslash{}\textbackslash{}{1, 2, \textbackslash{}\textbackslash{}dots, n-1\textbackslash{}\textbackslash{}\textbackslash{}}\textbackslash{}$. Then the alternating sum of the absolute values of the autocorrelations is bounded by:\textbackslash{}n\textbackslash{}n\textbackslash{}$\textbackslash{}$ \textbackslash{}\textbackslash{}sum\textbackslash{}_\textbackslash{}{k=1\textbackslash{}}\textbackslash{}textasciicircum{}\textbackslash{}{n-1\textbackslash{}} (-1)\textbackslash{}textasciicircum{}\textbackslash{}{k-1\textbackslash{}} |C\textbackslash{}_k(s)| \textbackslash{}\textbackslash{}le \textbackslash{}\textbackslash{}left\textbackslash{}\textbackslash{}lfloor \textbackslash{}\textbackslash{}frac\textbackslash{}{n\textbackslash{}}\textbackslash{}{2\textbackslash{}} \textbackslash{}\textbackslash{}right\textbackslash{}\textbackslash{}rfloor \textbackslash{}$\textbackslash{}$
\end{conjecture}

\begin{lstlisting}[caption={Lean 4 Sketch (v14 Galois)}]
import Mathlib.Data.Int.Basic
import Mathlib.Data.Fin.VecNotation
import Mathlib.Data.Finset.Sum
import Mathlib.Tactic.Linarith

open Finset BigOperators

-- A sequence of length n with elements in {-1, 1}
abbrev BinarySeq (n : ℕ) := {s : Fin n → ℤ // ∀ i, s i = 1 ∨ s i = -1}

-- Aperiodic autocorrelation C_k(s)
def autocorr {n : ℕ} (s : Fin n → ℤ) (k : ℕ) : ℤ := 
  if h : k > 0 ∧ k < n then
    let N := n - k
    -- sum over i from 0 to n-k-1
    ∑ i : Fin N, s ⟨i, by linarith [i.isLt, h.2]⟩ * s ⟨i + k, by linarith [i.isLt, h.2]⟩
  else
    0

theorem autocorrelation_signed_abs_upper_bound (n
\end{lstlisting}

\hrule\medskip

\subsection*{\texttt{calabi\textbackslash{}_yau\textbackslash{}_c5}}

\textbf{Domain:} special\textbackslash{}_functions \quad \textbf{Class:} 4 \quad \textbf{Verdict:} \textcolor{incompleteamber}{\textbf{INCOMPLETE}}

\textbf{Confidence:} 0.70 \quad \textbf{Domain Prior:} 0.25 \quad \textbf{Sorry gaps:} 5

\textcolor{sorryred}{\textbf{Zero-Sorry Guillotine Applied:} 5 \texttt{sorry} stubs detected. Verdict downgraded from VERIFIED to INCOMPLETE.}

\begin{conjecture}
\textbackslash{}\textbackslash{}int\textbackslash{}_0\textbackslash{}textasciicircum{}1 \textbackslash{}\textbackslash{}frac\textbackslash{}{\textbackslash{}\textbackslash{}mathrm\textbackslash{}{Li\textbackslash{}}\textbackslash{}_3(x)\textbackslash{}\textbackslash{}log(1-x)\textbackslash{}}\textbackslash{}{x\textbackslash{}}dx = \textbackslash{}\textbackslash{}zeta(2)\textbackslash{}\textbackslash{}zeta(3) - 3\textbackslash{}\textbackslash{}zeta(5)
\end{conjecture}

\begin{lstlisting}[caption={Lean 4 Sketch (v14 Galois)}]
import Mathlib.Analysis.SpecialFunctions.Log.Polylog
import Mathlib.Analysis.SpecialFunctions.Zeta.Basic
import Mathlib.Analysis.Calculus.Integration.Api

open Real Set Filter Topology

-- Define the integrand of the conjectural identity.
noncomputable def integrand (x : ℝ) : ℝ :=
  (polylog 3 x * log (1 - x)) / x

-- The formal statement of the conjecture.
theorem CalabiYauC5_Conjecture : 
    ∫ x in Ioc 0 1, integrand x = zeta 2 * zeta 3 - 3 * zeta 5 := by
  -- The proof proceeds by expanding log(1-x) and polylog 3 x into their series
  -- representations, integrating term-by-term, and evalu
\end{lstlisting}

\hrule\medskip

\subsection*{\texttt{knot\textbackslash{}_volume\textbackslash{}_7\textbackslash{}_2}}

\textbf{Domain:} discrete\textbackslash{}_geometry \quad \textbf{Class:} 3 \quad \textbf{Verdict:} \textcolor{incompleteamber}{\textbf{INCOMPLETE}}

\textbf{Confidence:} 0.70 \quad \textbf{Domain Prior:} 0.20 \quad \textbf{Sorry gaps:} 7

\textcolor{sorryred}{\textbf{Zero-Sorry Guillotine Applied:} 7 \texttt{sorry} stubs detected. Verdict downgraded from VERIFIED to INCOMPLETE.}

\begin{conjecture}
Let \textbackslash{}$K\textbackslash{}$ be the \textbackslash{}$7\textbackslash{}_2\textbackslash{}$ knot. Let \textbackslash{}$s(K) = 7\textbackslash{}$ be its stick number and let \textbackslash{}$\textbackslash{}\textbackslash{}mathcal\textbackslash{}{P\textbackslash{}}\textbackslash{}_n(K)\textbackslash{}$ denote the space of all \textbackslash{}$n\textbackslash{}$-vertex polygonal knots in \textbackslash{}$\textbackslash{}\textbackslash{}mathbb\textbackslash{}{R\textbackslash{}}\textbackslash{}textasciicircum{}3\textbackslash{}$ of type \textbackslash{}$K\textbackslash{}$. For a polygon \textbackslash{}$P \textbackslash{}\textbackslash{}in \textbackslash{}\textbackslash{}mathcal\textbackslash{}{P\textbackslash{}}\textbackslash{}_n(K)\textbackslash{}$, let \textbackslash{}$\textbackslash{}\textbackslash{}text\textbackslash{}{Vol\textbackslash{}}(\textbackslash{}\textbackslash{}text\textbackslash{}{conv\textbackslash{}}(P))\textbackslash{}$ be the Euclidean volume of the convex hull of its vertices and let \textbackslash{}$L(P)\textbackslash{}$ be its total length. Let \textbackslash{}$V\textbackslash{}_\textbackslash{}{\textbackslash{}\textbackslash{}text\textbackslash{}{hyp\textbackslash{}}\textbackslash{}}(K)\textbackslash{}$ be the hyperbolic volume of th
\end{conjecture}

\begin{lstlisting}[caption={Lean 4 Sketch (v14 Galois)}]
import Mathlib.Topology.KnotTheory
import Mathlib.Geometry.Euclidean.Basic
import Mathlib.Analysis.SpecialFunctions.Trigonometric.Arctan

-- We would need a formalization of polygonal knots and their properties
-- For now, we use stubs.

abbrev PolygonalKnot := List (EuclideanSpace ℝ (Fin 3))

-- A rigorous definition of knot type equivalence for polygonal knots
def IsKnotTypeOf (P1 P2 : PolygonalKnot) : Prop := sorry

-- The specific knot type 7_2
axiom k7_2_type : Type
axiom k7_2 : k7_2_type

-- A predicate for a polygon to be of a certain knot type
def IsOfKnotType (P : PolygonalKnot) (k : 
\end{lstlisting}

\hrule\medskip

\subsection*{\texttt{anderson\textbackslash{}_lyapunov\textbackslash{}_exponent}}

\textbf{Domain:} mathematical\textbackslash{}_physics \quad \textbf{Class:} 4 \quad \textbf{Verdict:} \textcolor{incompleteamber}{\textbf{INCOMPLETE}}

\textbf{Confidence:} 0.70 \quad \textbf{Domain Prior:} 0.15 \quad \textbf{Sorry gaps:} 4

\textcolor{sorryred}{\textbf{Zero-Sorry Guillotine Applied:} 4 \texttt{sorry} stubs detected. Verdict downgraded from VERIFIED to INCOMPLETE.}

\begin{conjecture}
Let \textbackslash{}$\textbackslash{}\textbackslash{}\textbackslash{}\textbackslash{}gamma(E, W)\textbackslash{}$ be the Lyapunov exponent for the one-dimensional Anderson-Bernoulli model on \textbackslash{}$\textbackslash{}\textbackslash{}\textbackslash{}\textbackslash{}mathbb\textbackslash{}{Z\textbackslash{}}\textbackslash{}$ with Hamiltonian \textbackslash{}$(H\textbackslash{}_\textbackslash{}\textbackslash{}\textbackslash{}\textbackslash{}omega \textbackslash{}\textbackslash{}\textbackslash{}\textbackslash{}psi)\textbackslash{}_n = \textbackslash{}\textbackslash{}\textbackslash{}\textbackslash{}psi\textbackslash{}_\textbackslash{}{n+1\textbackslash{}} + \textbackslash{}\textbackslash{}\textbackslash{}\textbackslash{}psi\textbackslash{}_\textbackslash{}{n-1\textbackslash{}} + W\textbackslash{}\textbackslash{}\textbackslash{}\textbackslash{}sigma\textbackslash{}_n \textbackslash{}\textbackslash{}\textbackslash{}\textbackslash{}psi\textbackslash{}_n\textbackslash{}$, where \textbackslash{}$\textbackslash{}\textbackslash{}\textbackslash{}\textbackslash{}sigma\textbackslash{}_n\textbackslash{}$ are i.i.d. random variables with \textbackslash{}$\textbackslash{}\textbackslash{}\textbackslash{}\textbackslash{}mathbb\textbackslash{}{P\textbackslash{}}(\textbackslash{}\textbackslash{}\textbackslash{}\textbackslash{}sigma\textbackslash{}_n = \textbackslash{}\textbackslash{}\textbackslash{}\textbackslash{}pm 1) = 1/2\textbackslash{}$. Let \textbackslash{}$T\textbackslash{}_\textbackslash{}{\textbackslash{}\textbackslash{}\textbackslash{}\textbackslash{}pm\textbackslash{}}(E,W) = \textbackslash{}\textbackslash{}\textbackslash{}\textbackslash{}begin\textbackslash{}{pmatrix\textbackslash{}} E \textbackslash{}\textbackslash{}\textbackslash{}\textbackslash{}mp W \textbackslash{}& -1 \textbackslash{}\textbackslash{}\textbackslash{}\textbackslash{}\textbackslash{}\textbackslash{}\textbackslash{}\textbackslash{} 1 \textbackslash{}& 0 \textbackslash{}\textbackslash{}\textbackslash{}\textbackslash{}end\textbackslash{}{pma
\end{conjecture}

\begin{lstlisting}[caption={Lean 4 Sketch (v14 Galois)}]
import Mathlib.Analysis.Analytic.Basic\
import Mathlib.LinearAlgebra.Matrix.SpecialLinearGroup\
import Mathlib.Probability.MeasureTheory.Measure.ProbabilityMeasure\
\
open Matrix Set MeasureTheory\
\
-- We postulate the existence of the Lyapunov exponent and its properties,\
-- as its full formalization is a major project.\
noncomputable axiom lyapunovExponent (E W : ℝ) : ℝ\
\
-- We define the transfer matrices for the 1D Anderson-Bernoulli model.\
noncomputable def transferMatrix (E W : ℝ) (s : Bool) : SL(2, ℝ) :=\
  let V := if s then W else -W\
  let M : Matrix (Fin 2) (Fin 2) ℝ := !![E - V
\end{lstlisting}

\hrule\medskip

\subsection*{\texttt{quartic\textbackslash{}_oscillator\textbackslash{}_lambda}}

\textbf{Domain:} spectral\textbackslash{}_theory \quad \textbf{Class:} 3 \quad \textbf{Verdict:} \textcolor{incompleteamber}{\textbf{INCOMPLETE}}

\textbf{Confidence:} 0.70 \quad \textbf{Domain Prior:} 0.20 \quad \textbf{Sorry gaps:} 4

\textcolor{sorryred}{\textbf{Zero-Sorry Guillotine Applied:} 4 \texttt{sorry} stubs detected. Verdict downgraded from VERIFIED to INCOMPLETE.}

\begin{conjecture}
Let \textbackslash{}$H = -\textbackslash{}\textbackslash{}frac\textbackslash{}{d\textbackslash{}textasciicircum{}2\textbackslash{}}\textbackslash{}{dx\textbackslash{}textasciicircum{}2\textbackslash{}} + x\textbackslash{}textasciicircum{}4\textbackslash{}$ be the Schrödinger operator on \textbackslash{}$L\textbackslash{}textasciicircum{}2(\textbackslash{}\textbackslash{}mathbb\textbackslash{}{R\textbackslash{}})\textbackslash{}$. Let its discrete, positive eigenvalues be ordered as \textbackslash{}$0 < \textbackslash{}\textbackslash{}lambda\textbackslash{}_0 < \textbackslash{}\textbackslash{}lambda\textbackslash{}_1 < \textbackslash{}\textbackslash{}lambda\textbackslash{}_2 < \textbackslash{}\textbackslash{}dots\textbackslash{}$. Then the alternating sum of the reciprocals of these eigenvalues evaluates to a specific constant related to the period of the classical quartic oscillator:\textbackslash{}n\textbackslash{}n\textbackslash{}$\textbackslash{}$\textbackslash{}\textbackslash{}sum\textbackslash{}_\textbackslash{}{n=0\textbackslash{}}\textbackslash{}textasciicircum{}\textbackslash{}\textbackslash{}infty \textbackslash{}\textbackslash{}frac\textbackslash{}{(-1)\textbackslash{}textasciicircum{}n\textbackslash{}}\textbackslash{}{\textbackslash{}\textbackslash{}lambda\textbackslash{}_n\textbackslash{}}
\end{conjecture}

\begin{lstlisting}[caption={Lean 4 Sketch (v14 Galois)}]
import Mathlib.Analysis.SpecialFunctions.Gamma.Basic
import Mathlib.Analysis.InnerProductSpace.Spectrum

open Real Filter Topology

-- This conjecture concerns the eigenvalues of the Schrödinger operator with a quartic potential.
-- The full definition of this unbounded operator and the proof of its spectral properties
-- are beyond the current scope of Mathlib. We axiomatically assume the existence of its
-- ordered sequence of eigenvalues with standard properties.
axiom quartic_eigenvalues : ℕ → ℝ
axiom quartic_eigenvalues_pos (n : ℕ) : 0 < quartic_eigenvalues n
axiom quartic_eigenvalues_str
\end{lstlisting}

\hrule\medskip

\subsection*{\texttt{spheroidal\textbackslash{}_eigenvalue\textbackslash{}_lambda\textbackslash{}_m0}}

\textbf{Domain:} spectral\textbackslash{}_theory \quad \textbf{Class:} 3 \quad \textbf{Verdict:} \textcolor{incompleteamber}{\textbf{INCOMPLETE}}

\textbf{Confidence:} 0.70 \quad \textbf{Domain Prior:} 0.20 \quad \textbf{Sorry gaps:} 1

\textcolor{sorryred}{\textbf{Zero-Sorry Guillotine Applied:} 1 \texttt{sorry} stubs detected. Verdict downgraded from VERIFIED to INCOMPLETE.}

\begin{conjecture}
Let \textbackslash{}$\textbackslash{}\textbackslash{}lambda\textbackslash{}_\textbackslash{}{0,n\textbackslash{}}(c)\textbackslash{}$ for \textbackslash{}$n \textbackslash{}\textbackslash{}in \textbackslash{}\textbackslash{}mathbb\textbackslash{}{N\textbackslash{}}\textbackslash{}$ be the sequence of eigenvalues of the angular prolate spheroidal differential equation, indexed in increasing order, for a fixed parameter \textbackslash{}$c \textbackslash{}\textbackslash{}in \textbackslash{}\textbackslash{}mathbb\textbackslash{}{R\textbackslash{}}\textbackslash{}_\textbackslash{}{>0\textbackslash{}}\textbackslash{}$. The equation is given by \textbackslash{}$\textbackslash{}\textbackslash{}frac\textbackslash{}{d\textbackslash{}}\textbackslash{}{dx\textbackslash{}}\textbackslash{}\textbackslash{}left((1-x\textbackslash{}textasciicircum{}2)\textbackslash{}\textbackslash{}frac\textbackslash{}{dS\textbackslash{}}\textbackslash{}{dx\textbackslash{}}\textbackslash{}\textbackslash{}right) + \textbackslash{}\textbackslash{}left(\textbackslash{}\textbackslash{}lambda - c\textbackslash{}textasciicircum{}2 x\textbackslash{}textasciicircum{}2\textbackslash{}\textbackslash{}right)S = 0\textbackslash{}$ on the interval \textbackslash{}$x \textbackslash{}\textbackslash{}in [-1, 1]\textbackslash{}$, with the condition that the solut
\end{conjecture}

\begin{lstlisting}[caption={Lean 4 Sketch (v14 Galois)}]
import Mathlib.Analysis.SpecialFunctions.Legendre
import Mathlib.Data.Real.Basic

/-!
# Convexity of Spheroidal Eigenvalues

This file states a conjecture about the eigenvalues of the prolate spheroidal wave equation
for azimuthal quantum number m=0. We postulate the existence of these eigenvalues as a
function `spheroidalEigenvalue` and state some of its known properties.

A full formalization would require defining the spheroidal differential operator on a
weighted Sobolev space and proving the existence of a discrete, ordered spectrum via the
spectral theory of self-adjoint operators.
-/

/
\end{lstlisting}

\hrule\medskip

\subsection*{\texttt{nested\textbackslash{}_radical\textbackslash{}_kasner}}

\textbf{Domain:} number\textbackslash{}_theory \quad \textbf{Class:} 3 \quad \textbf{Verdict:} \textcolor{incompleteamber}{\textbf{INCOMPLETE}}

\textbf{Confidence:} 0.70 \quad \textbf{Domain Prior:} 0.05 \quad \textbf{Sorry gaps:} 5

\textcolor{sorryred}{\textbf{Zero-Sorry Guillotine Applied:} 5 \texttt{sorry} stubs detected. Verdict downgraded from VERIFIED to INCOMPLETE.}

\begin{conjecture}
Let \textbackslash{}$(a\textbackslash{}_n)\textbackslash{}_\textbackslash{}{n \textbackslash{}\textbackslash{}ge 1\textbackslash{}}\textbackslash{}$ be a sequence of positive integers defined by a polynomial \textbackslash{}$Q(n) \textbackslash{}\textbackslash{}in \textbackslash{}\textbackslash{}mathbb\textbackslash{}{Z\textbackslash{}}[n]\textbackslash{}$ of degree \textbackslash{}$2d\textbackslash{}$, i.e., \textbackslash{}$a\textbackslash{}_n = Q(n)\textbackslash{}$ for all \textbackslash{}$n \textbackslash{}\textbackslash{}ge 1\textbackslash{}$. The nested radical \textbackslash{}$X = \textbackslash{}\textbackslash{}sqrt\textbackslash{}{a\textbackslash{}_1 + \textbackslash{}\textbackslash{}sqrt\textbackslash{}{a\textbackslash{}_2 + \textbackslash{}\textbackslash{}sqrt\textbackslash{}{a\textbackslash{}_3 + \textbackslash{}\textbackslash{}dots\textbackslash{}}\textbackslash{}}\textbackslash{}}\textbackslash{}$ converges to an integer \textbackslash{}$N\textbackslash{}$ if and only if there exists a polynomial \textbackslash{}$P(n) \textbackslash{}\textbackslash{}in \textbackslash{}\textbackslash{}mathbb\textbackslash{}{Q\textbackslash{}}[n]\textbackslash{}$ of degree \textbackslash{}$d\textbackslash{}$ such that:\textbackslash{}n\textbackslash{}n1.  \textbackslash{}$Q(n) = P(n)\textbackslash{}textasciicircum{}2 - P(n+1)\textbackslash{}$ for all \textbackslash{}$
\end{conjecture}

\begin{lstlisting}[caption={Lean 4 Sketch (v14 Galois)}]
import Mathlib.Data.Polynomial.Basic
import Mathlib.Analysis.SpecificLimits.Basic
import Mathlib.Data.Rat.Defs

open Polynomial Filter

-- We assume the existence of a formal definition for the value of a nested radical,
-- which involves proving the convergence of the sequence of finite approximations.
-- `h_conv` would be a proof that the sequence converges.
def NestedRadical (a : ℕ → ℝ) (h_conv : ∃ l, Tendsto (fun n ↦ sorry) atTop (𝓝 l)) : ℝ := sorry

-- The conjecture statement
theorem nested_radical_kasner_conjecture
  (Q : Polynomial ℤ) (h_pos : ∀ n : ℕ, n > 0 → 0 < (Q.eval (n : ℤ))) :
 
\end{lstlisting}

\hrule\medskip

\subsection*{\texttt{mrb\textbackslash{}_constant}}

\textbf{Domain:} number\textbackslash{}_theory \quad \textbf{Class:} 3 \quad \textbf{Verdict:} \textcolor{refutedred}{\textbf{REFUTED}}

\textbf{Confidence:} 0.57 \quad \textbf{Domain Prior:} 0.05 \quad \textbf{Sorry gaps:} 3

\begin{conjecture}
Let \textbackslash{}$\textbackslash{}\textbackslash{}mathcal\textbackslash{}{P\textbackslash{}}\textbackslash{}$ be the ring of periods, defined as the set of all real numbers that can be expressed as the value of an absolutely convergent integral of a rational function with rational coefficients over a semi-algebraic set in \textbackslash{}$\textbackslash{}\textbackslash{}mathbb\textbackslash{}{R\textbackslash{}}\textbackslash{}textasciicircum{}n\textbackslash{}$ defined by polynomial inequalities with rational coefficients. The MRB constant, \textbackslash{}$C\textbackslash{}_\textbackslash{}{MRB\textbackslash{}} = \textbackslash{}\textbackslash{}int\textbackslash{}_0\textbackslash{}textasciicircum{}1 x\textbackslash{}textasciicircum{}\textbackslash{}{-x\textbackslash{}} dx = \textbackslash{}\textbackslash{}sum\textbackslash{}_\textbackslash{}{n=1\textbackslash{}}\textbackslash{}textasciicircum{}\textbackslash{}\textbackslash{}infty n\textbackslash{}textasciicircum{}\textbackslash{}{-n\textbackslash{}}\textbackslash{}$, is not a
\end{conjecture}

\begin{lstlisting}[caption={Lean 4 Sketch (v14 Galois)}]
import Mathlib.Analysis.SpecialFunctions.Pow.Real
import Mathlib.Analysis.SpecialFunctions.Integrals
import Mathlib.NumberTheory.Liouville.Basic

open Real Set MeasureTheory Filter

-- We define the MRB constant via its integral representation.
-- The integrand x ↦ x⁻ˣ is continuous on (0, 1], and its limit at 0 is 1.
-- We define a function equal to x⁻ˣ on (0, 1] and 1 at 0, which is continuous on [0, 1].
noncomputable def mrbIntegrand' (x : ℝ) : ℝ := if x > 0 then x ^ (-x) else 1

-- A proof of continuity on the closed interval [0,1].
lemma mrbIntegrand'_continuous_on : ContinuousOn mrbInteg
\end{lstlisting}

\hrule\medskip

\subsection*{\texttt{torsional\textbackslash{}_rigidity\textbackslash{}_square}}

\textbf{Domain:} special\textbackslash{}_functions \quad \textbf{Class:} 3 \quad \textbf{Verdict:} \textcolor{incompleteamber}{\textbf{INCOMPLETE}}

\textbf{Confidence:} 0.70 \quad \textbf{Domain Prior:} 0.25 \quad \textbf{Sorry gaps:} 4

\textcolor{sorryred}{\textbf{Zero-Sorry Guillotine Applied:} 4 \texttt{sorry} stubs detected. Verdict downgraded from VERIFIED to INCOMPLETE.}

\begin{conjecture}
Let \textbackslash{}$S = \textbackslash{}\textbackslash{}sum\textbackslash{}_\textbackslash{}{n,m \textbackslash{}\textbackslash{}in 2\textbackslash{}\textbackslash{}mathbb\textbackslash{}{N\textbackslash{}}-1\textbackslash{}} \textbackslash{}\textbackslash{}frac\textbackslash{}{1\textbackslash{}}\textbackslash{}{n\textbackslash{}textasciicircum{}2 m\textbackslash{}textasciicircum{}2 (n\textbackslash{}textasciicircum{}2+m\textbackslash{}textasciicircum{}2)\textbackslash{}}\textbackslash{}$. The value of this sum is given by the identity: \textbackslash{}n\textbackslash{}$\textbackslash{}$ S = \textbackslash{}\textbackslash{}frac\textbackslash{}{\textbackslash{}\textbackslash{}pi\textbackslash{}textasciicircum{}6\textbackslash{}}\textbackslash{}{2048\textbackslash{}} \textbackslash{}\textbackslash{}left( \textbackslash{}\textbackslash{}frac\textbackslash{}{16\textbackslash{}}\textbackslash{}{3\textbackslash{}} - \textbackslash{}\textbackslash{}frac\textbackslash{}{1024\textbackslash{}}\textbackslash{}{\textbackslash{}\textbackslash{}pi\textbackslash{}textasciicircum{}5\textbackslash{}} \textbackslash{}\textbackslash{}sum\textbackslash{}_\textbackslash{}{k=1, k \textbackslash{}\textbackslash{}text\textbackslash{}{ odd\textbackslash{}}\textbackslash{}}\textbackslash{}}\textbackslash{}textasciicircum{}\textbackslash{}{\textbackslash{}\textbackslash{}infty\textbackslash{}} \textbackslash{}\textbackslash{}frac\textbackslash{}{\textbackslash{}\textbackslash{}tanh(k\textbackslash{}\textbackslash{}pi/2)\textbackslash{}}\textbackslash{}{k\textbackslash{}textasciicircum{}5\textbackslash{}} \textbackslash{}\textbackslash{}right) \textbackslash{}$\textbackslash{}$
\end{conjecture}

\begin{lstlisting}[caption={Lean 4 Sketch (v14 Galois)}]
import Mathlib.Analysis.SpecialFunctions.Trigonometric.Series
import Mathlib.Analysis.Fourier.FourierSeries
import Mathlib.Analysis.Calculus.ParametricIntegral

open Nat Real BigOperators Filter Topology

/- 
  Let S be the double summation over odd natural numbers n, m of 1 / (n^2 * m^2 * (n^2 + m^2)).
  This conjecture states that S can be expressed in terms of a single, faster-converging series
  involving the hyperbolic tangent function. This identity is motivated by equating two distinct
  series representations for the torsional rigidity of a square domain. 
-/
theorem torsional_rigidity
\end{lstlisting}

\hrule\medskip

\subsection*{\texttt{schur\textbackslash{}_6}}

\textbf{Domain:} combinatorics \quad \textbf{Class:} 4 \quad \textbf{Verdict:} \textcolor{incompleteamber}{\textbf{INCOMPLETE}}

\textbf{Confidence:} 0.70 \quad \textbf{Domain Prior:} 0.35 \quad \textbf{Sorry gaps:} 3

\textcolor{sorryred}{\textbf{Zero-Sorry Guillotine Applied:} 3 \texttt{sorry} stubs detected. Verdict downgraded from VERIFIED to INCOMPLETE.}

\begin{conjecture}
Let \textbackslash{}$N\textbackslash{}_6(c)\textbackslash{}$ be the least integer \textbackslash{}$N\textbackslash{}$ such that for any coloring \textbackslash{}$f: \textbackslash{}\textbackslash{}\textbackslash{}{1, \textbackslash{}\textbackslash{}dots, N\textbackslash{}\textbackslash{}\textbackslash{}} \textbackslash{}\textbackslash{}to \textbackslash{}\textbackslash{}\textbackslash{}{1, \textbackslash{}\textbackslash{}dots, c\textbackslash{}\textbackslash{}\textbackslash{}}\textbackslash{}$, there exist \textbackslash{}$x\textbackslash{}_1, \textbackslash{}\textbackslash{}dots, x\textbackslash{}_6 \textbackslash{}\textbackslash{}in \textbackslash{}\textbackslash{}\textbackslash{}{1, \textbackslash{}\textbackslash{}dots, N\textbackslash{}\textbackslash{}\textbackslash{}}\textbackslash{}$ satisfying \textbackslash{}$f(x\textbackslash{}_1) = f(x\textbackslash{}_2) = \textbackslash{}\textbackslash{}dots = f(x\textbackslash{}_6)\textbackslash{}$ and \textbackslash{}$x\textbackslash{}_1+x\textbackslash{}_2+x\textbackslash{}_3=x\textbackslash{}_4+x\textbackslash{}_5+x\textbackslash{}_6\textbackslash{}$. Then the asymptotic growth rate of \textbackslash{}$N\textbackslash{}_6(c)\textbackslash{}$ is given by \textbackslash{}$N\textbackslash{}_6(c) = \textbackslash{}\textbackslash{}Theta(c\textbackslash{}textasciicircum{}\textbackslash{}{3/2\textbackslash{}})\textbackslash{}$.
\end{conjecture}

\begin{lstlisting}[caption={Lean 4 Sketch (v14 Galois)}]
import Mathlib.Analysis.Asymptotics.Asymptotics
import Mathlib.Data.Finset.Card
import Mathlib.Data.Nat.Basic
import Mathlib.Tactic

open Finset Asymptotics Filter
open scoped BigOperators Nat

/-- A finite set of natural numbers `S` has a Schur-6 solution if there exist
    elements (not necessarily distinct) `x, y, z, a, b, c` in `S` such that `x+y+z = a+b+c`. -/
def HasSchur6Solution (S : Finset ℕ) : Prop := 
  ∃ x y z a b c, x ∈ S ∧ y ∈ S ∧ z ∈ S ∧ a ∈ S ∧ b ∈ S ∧ c ∈ S ∧ x + y + z = a + b + c

/-- The Schur-6 number for `c` colors, `Schur6Number c`, is the smallest `n` such that any
    `
\end{lstlisting}

\hrule\medskip

\subsection*{\texttt{diff\textbackslash{}_basis\textbackslash{}_optimal\textbackslash{}_10000}}

\textbf{Domain:} combinatorics \quad \textbf{Class:} 3 \quad \textbf{Verdict:} \textcolor{refutedred}{\textbf{REFUTED}}

\textbf{Confidence:} 0.57 \quad \textbf{Domain Prior:} 0.35 \quad \textbf{Sorry gaps:} 1

\begin{conjecture}
Let \textbackslash{}$S \textbackslash{}\textbackslash{}subset \textbackslash{}\textbackslash{}mathbb\textbackslash{}{Z\textbackslash{}}\textbackslash{}$ be a finite set. Define the set of positive differences as \textbackslash{}$D(S) := \textbackslash{}\textbackslash{}\textbackslash{}{ s\textbackslash{}_i - s\textbackslash{}_j \textbackslash{}\textbackslash{}mid s\textbackslash{}_i, s\textbackslash{}_j \textbackslash{}\textbackslash{}in S, s\textbackslash{}_i > s\textbackslash{}_j \textbackslash{}\textbackslash{}\textbackslash{}}\textbackslash{}$. Define the contiguous range covered by these differences as \textbackslash{}$n(S) := \textbackslash{}\textbackslash{}sup \textbackslash{}\textbackslash{}\textbackslash{}{ n \textbackslash{}\textbackslash{}in \textbackslash{}\textbackslash{}mathbb\textbackslash{}{N\textbackslash{}} \textbackslash{}\textbackslash{}cup \textbackslash{}\textbackslash{}\textbackslash{}{0\textbackslash{}\textbackslash{}\textbackslash{}} \textbackslash{}\textbackslash{}mid \textbackslash{}\textbackslash{}\textbackslash{}{1, 2, \textbackslash{}\textbackslash{}dots, n\textbackslash{}\textbackslash{}\textbackslash{}} \textbackslash{}\textbackslash{}subseteq D(S) \textbackslash{}\textbackslash{}\textbackslash{}}\textbackslash{}$. A set \textbackslash{}$S\textbackslash{}$ is called an optimal difference basis if for any set \textbackslash{}$S'\textbackslash{}$ with \textbackslash{}$|S'| < |S|\textbackslash{}$, w
\end{conjecture}

\begin{lstlisting}[caption={Lean 4 Sketch (v14 Galois)}]
import Mathlib.Data.Finset.Pointwise
import Mathlib.Data.Finset.Card
import Mathlib.Tactic

-- Let S be a finite set of integers
variable (S : Finset ℤ)

-- The set of positive differences of S
def positive_differences (S : Finset ℤ) : Finset ℤ := 
  Finset.image (fun p : ℤ × ℤ ↦ p.2 - p.1) (S.offDiag.filter (fun p ↦ p.1 < p.2))

-- The length of the initial segment of integers covered by the differences of S
-- We search up to a loose bound Nat.choose S.card 2, the max possible number of positive diffs
def n (S : Finset ℤ) : ℕ :=
  if h : S.card < 2 then 0
  else Nat.findGreatest 
    (fun m 
\end{lstlisting}

\hrule\medskip

\subsection*{\texttt{general\textbackslash{}_diff\textbackslash{}_basis\textbackslash{}_algo}}

\textbf{Domain:} combinatorics \quad \textbf{Class:} 3 \quad \textbf{Verdict:} \textcolor{refutedred}{\textbf{REFUTED}}

\textbf{Confidence:} 0.62 \quad \textbf{Domain Prior:} 0.35 \quad \textbf{Sorry gaps:} 3

\begin{conjecture}
Let \textbackslash{}$(Q, *)\textbackslash{}$ be a finite quasigroup of order \textbackslash{}$n\textbackslash{}$. Let the right division operator be defined by \textbackslash{}$x/y = z\textbackslash{}$ if and only if \textbackslash{}$z*y=x\textbackslash{}$. A set \textbackslash{}$B \textbackslash{}\textbackslash{}\textbackslash{}\textbackslash{}subseteq Q\textbackslash{}$ is a right quotient basis if the set of quotients \textbackslash{}$\textbackslash{}\textbackslash{}\textbackslash{}{x/y \textbackslash{}\textbackslash{}\textbackslash{}\textbackslash{}mid x, y \textbackslash{}\textbackslash{}\textbackslash{}\textbackslash{}in B\textbackslash{}\textbackslash{}\textbackslash{}}\textbackslash{}$ is equal to \textbackslash{}$Q\textbackslash{}$. Let \textbackslash{}$\textbackslash{}\textbackslash{}\textbackslash{}\textbackslash{}mu\textbackslash{}_r(Q)\textbackslash{}$ denote the minimum size of a right quotient basis for \textbackslash{}$Q\textbackslash{}$. Let \textbackslash{}$\textbackslash{}\textbackslash{}\textbackslash{}\textbackslash{}text\textbackslash{}{Mlt\textbackslash{}}\textbackslash{}_r(Q)\textbackslash{}$ be the right multiplication group of \textbackslash{}$Q\textbackslash{}$, 
\end{conjecture}

\begin{lstlisting}[caption={Lean 4 Sketch (v14 Galois)}]
import Mathlib.Algebra.Group.Defs\
import Mathlib.GroupTheory.Solvable\
import Mathlib.GroupTheory.Subgroup.Basic\
import Mathlib.Data.Set.Finite\
import Mathlib.Data.Nat.Log\
import Mathlib.Tactic\
\
universe u\
\
variable {Q : Type u} [Quasigroup Q] [Fintype Q]\
\
/-- The set of elements representable as a right quotient of elements from a set B. -/\
def right_quotient_set (B : Set Q) : Set Q := \
  {q | ∃ x ∈ B, ∃ y ∈ B, q = x / y}\
\
/-- A set B is a right quotient basis if its right quotient set is the entire quasigroup. -/\
def is_right_quotient_basis (B : Set Q) : Prop := \
  right_quot
\end{lstlisting}

\hrule\medskip

\subsection*{\texttt{merit\textbackslash{}_factor\textbackslash{}_6\textbackslash{}_5}}

\textbf{Domain:} coding\textbackslash{}_theory \quad \textbf{Class:} 3 \quad \textbf{Verdict:} \textcolor{incompleteamber}{\textbf{INCOMPLETE}}

\textbf{Confidence:} 0.70 \quad \textbf{Domain Prior:} 0.30 \quad \textbf{Sorry gaps:} 1

\textcolor{sorryred}{\textbf{Zero-Sorry Guillotine Applied:} 1 \texttt{sorry} stubs detected. Verdict downgraded from VERIFIED to INCOMPLETE.}

\begin{conjecture}
Let \textbackslash{}$p\textbackslash{}$ be a prime such that \textbackslash{}$p \textbackslash{}\textbackslash{}equiv 1 \textbackslash{}\textbackslash{}pmod 4\textbackslash{}$. Let \textbackslash{}$L\textbackslash{}_p = ( (\textbackslash{}\textbackslash{}frac\textbackslash{}{k\textbackslash{}}\textbackslash{}{p\textbackslash{}}) )\textbackslash{}_\textbackslash{}{k=1, \textbackslash{}\textbackslash{}dots, p-1\textbackslash{}}\textbackslash{}$ be the Legendre sequence of length \textbackslash{}$n=p-1\textbackslash{}$. Let \textbackslash{}$\textbackslash{}\textbackslash{}mathcal\textbackslash{}{A\textbackslash{}}\textbackslash{}_p\textbackslash{}$ be the set of all \textbackslash{}$n\textbackslash{}$ cyclic shifts of \textbackslash{}$L\textbackslash{}_p\textbackslash{}$. For a sequence \textbackslash{}$A\textbackslash{}$ of length \textbackslash{}$n\textbackslash{}$, define the 'energy' as \textbackslash{}$E(A) = 2 \textbackslash{}\textbackslash{}sum\textbackslash{}_\textbackslash{}{k=1\textbackslash{}}\textbackslash{}textasciicircum{}\textbackslash{}{n-1\textbackslash{}} C\textbackslash{}_k(A)\textbackslash{}textasciicircum{}2\textbackslash{}$, where \textbackslash{}$C\textbackslash{}_k(A)\textbackslash{}$ are the aperiodic autocorrelation coefficients. Let \textbackslash{}$E\textbackslash{}_\textbackslash{}{\textbackslash{}\textbackslash{}text\textbackslash{}{opt\textbackslash{}}\textbackslash{}}(p) = \textbackslash{}
\end{conjecture}

\begin{lstlisting}[caption={Lean 4 Sketch (v14 Galois)}]
import Mathlib.NumberTheory.LegendreSymbol
import Mathlib.Analysis.Asymptotics.Asymptotics
import Mathlib.Data.Real.Basic
import Mathlib.Data.Fin.VecNotation
import Mathlib.Tactic
import Mathlib.Analysis.SpecialFunctions.Pow.Real

open Filter Real Nat Finset

noncomputable section

-- Define binary sequences as maps to {-1, 1}
abbrev BinarySeq (n : ℕ) := Fin n → ℤ

-- Aperiodic autocorrelation function
def aperiodicCorr (A : BinarySeq n) (k : ℕ) : ℤ :=
  if h : k < n ∧ k > 0 then
    ∑ i in range (n - k), A ⟨i, lt_of_lt_of_le i.2 (Nat.sub_le n k)⟩ * A ⟨i + k, Nat.add_lt_of_lt_of_le i.2 (Nat.su
\end{lstlisting}

\hrule\medskip

\subsection*{\texttt{parametric\textbackslash{}_spherical\textbackslash{}_codes}}

\textbf{Domain:} coding\textbackslash{}_theory \quad \textbf{Class:} 4 \quad \textbf{Verdict:} \textcolor{incompleteamber}{\textbf{INCOMPLETE}}

\textbf{Confidence:} 0.70 \quad \textbf{Domain Prior:} 0.30 \quad \textbf{Sorry gaps:} 1

\textcolor{sorryred}{\textbf{Zero-Sorry Guillotine Applied:} 1 \texttt{sorry} stubs detected. Verdict downgraded from VERIFIED to INCOMPLETE.}

\begin{conjecture}
Let \textbackslash{}$C \textbackslash{}\textbackslash{}subset S\textbackslash{}textasciicircum{}\textbackslash{}{n-1\textbackslash{}} \textbackslash{}\textbackslash{}subset \textbackslash{}\textbackslash{}mathbb\textbackslash{}{R\textbackslash{}}\textbackslash{}textasciicircum{}n\textbackslash{}$ be a finite set of points (a spherical code). If \textbackslash{}$C\textbackslash{}$ is DGS-tight, meaning it achieves the Delsarte-Goethals-Seidel linear programming bound for some auxiliary polynomial \textbackslash{}$f(t)\textbackslash{}$ derived from Gegenbauer polynomials, then the set of all inner products between distinct points of \textbackslash{}$C\textbackslash{}$, \textbackslash{}$A = \textbackslash{}\textbackslash{}\textbackslash{}{ \textbackslash{}\textbackslash{}langle x, y \textbackslash{}\textbackslash{}rangle \textbackslash{}\textbackslash{}mid x, y \textbackslash{}\textbackslash{}in C, x \textbackslash{}\textbackslash{}neq y \textbackslash{}\textbackslash{}\textbackslash{}}\textbackslash{}$, mus
\end{conjecture}

\begin{lstlisting}[caption={Lean 4 Sketch (v14 Galois)}]
import Mathlib.Analysis.InnerProductSpace.Basic
import Mathlib.Data.Real.Basic
import Mathlib.Data.Rat.Cast
import Mathlib.Topology.MetricSpace.Sphere
import Mathlib.Algebra.Polynomial.Roots

open InnerProductSpace Finset Real

-- Let E be a finite-dimensional real inner product space.
variable {E : Type*} [NormedAddCommGroup E] [InnerProductSpace ℝ E] [FiniteDimensional ℝ E]

/-- A spherical code is a finite set of points on the unit sphere. -/
def SphericalCode := { C : Finset E // ∀ x ∈ C, ‖x‖ = 1 }

/-- The set of inner products between distinct points of a code. -/
def innerProducts (C : 
\end{lstlisting}

\hrule\medskip

\subsection*{\texttt{bklc\textbackslash{}_68\textbackslash{}_15}}

\textbf{Domain:} coding\textbackslash{}_theory \quad \textbf{Class:} 3 \quad \textbf{Verdict:} \textcolor{incompleteamber}{\textbf{INCOMPLETE}}

\textbf{Confidence:} 0.70 \quad \textbf{Domain Prior:} 0.30 \quad \textbf{Sorry gaps:} 3

\textcolor{sorryred}{\textbf{Zero-Sorry Guillotine Applied:} 3 \texttt{sorry} stubs detected. Verdict downgraded from VERIFIED to INCOMPLETE.}

\begin{conjecture}
Let \textbackslash{}$C\textbackslash{}_\textbackslash{}{\textbackslash{}\textbackslash{}\textbackslash{}\textbackslash{}text\textbackslash{}{GRS\textbackslash{}}\textbackslash{}}\textbackslash{}$ be a Generalized Reed-Solomon code over the finite field \textbackslash{}$\textbackslash{}\textbackslash{}\textbackslash{}\textbackslash{}mathbb\textbackslash{}{F\textbackslash{}}\textbackslash{}_\textbackslash{}{16\textbackslash{}}\textbackslash{}$ with parameters \textbackslash{}$[n,k,d] = [17, 3, 15]\textbackslash{}$. Let \textbackslash{}$\textbackslash{}\textbackslash{}\textbackslash{}\textbackslash{}phi: \textbackslash{}\textbackslash{}\textbackslash{}\textbackslash{}mathbb\textbackslash{}{F\textbackslash{}}\textbackslash{}_\textbackslash{}{16\textbackslash{}} \textbackslash{}\textbackslash{}\textbackslash{}\textbackslash{}to (\textbackslash{}\textbackslash{}\textbackslash{}\textbackslash{}mathbb\textbackslash{}{F\textbackslash{}}\textbackslash{}_2)\textbackslash{}textasciicircum{}4\textbackslash{}$ be an isomorphism of \textbackslash{}$\textbackslash{}\textbackslash{}\textbackslash{}\textbackslash{}mathbb\textbackslash{}{F\textbackslash{}}\textbackslash{}_2\textbackslash{}$-vector spaces. The binary code \textbackslash{}$C\textbackslash{}_\textbackslash{}{\textbackslash{}\textbackslash{}\textbackslash{}\textbackslash{}text\textbackslash{}{bin\textbackslash{}}\textbackslash{}} = \textbackslash{}\textbackslash{}\textbackslash{}\textbackslash{}\textbackslash{}{ (\textbackslash{}\textbackslash{}\textbackslash{}\textbackslash{}phi(c\textbackslash{}_0), \textbackslash{}\textbackslash{}\textbackslash{}\textbackslash{}dots, \textbackslash{}\textbackslash{}phi(c\textbackslash{}_\textbackslash{}{16\textbackslash{}})) \textbackslash{}\textbackslash{}\textbackslash{}\textbackslash{}mid (c\textbackslash{}_0, \textbackslash{}\textbackslash{}\textbackslash{}\textbackslash{}dots, c\textbackslash{}_\textbackslash{}{16\textbackslash{}}) \textbackslash{}\textbackslash{}\textbackslash{}\textbackslash{}in C\textbackslash{}_\textbackslash{}{\textbackslash{}\textbackslash{}\textbackslash{}\textbackslash{}text\textbackslash{}{GRS\textbackslash{}}\textbackslash{}} \textbackslash{}\textbackslash{}
\end{conjecture}

\begin{lstlisting}[caption={Lean 4 Sketch (v14 Galois)}]
import Mathlib.Data.Fintype.Card\
import Mathlib.FieldTheory.Finite.Basic\
import Mathlib.LinearAlgebra.Basis\
import Mathlib.Code.LinearCode\
\
-- Conjectured setup\
namespace Bklc6815\
\
-- Let F be the finite field with 16 elements.\
variable (F : Type) [Field F] [Fintype F] (hF : Fintype.card F = 16)\
\
-- A GRS code is a type of evaluation code.\
-- We postulate its existence and key properties as it is not yet defined in Mathlib.\
-- A full proof would require constructing evaluation codes.\
-- Parameters: n=17, k=3, over F\
\
-- We model the code as a subspace of (Fin 17 → F)\
variable 
\end{lstlisting}

\hrule\medskip

\subsection*{\texttt{lattice\textbackslash{}_packing\textbackslash{}_dim10}}

\textbf{Domain:} discrete\textbackslash{}_geometry \quad \textbf{Class:} 3 \quad \textbf{Verdict:} \textcolor{incompleteamber}{\textbf{INCOMPLETE}}

\textbf{Confidence:} 0.70 \quad \textbf{Domain Prior:} 0.20 \quad \textbf{Sorry gaps:} 1

\textcolor{sorryred}{\textbf{Zero-Sorry Guillotine Applied:} 1 \texttt{sorry} stubs detected. Verdict downgraded from VERIFIED to INCOMPLETE.}

\begin{conjecture}
Let \textbackslash{}$\textbackslash{}\textbackslash{}Lambda \textbackslash{}\textbackslash{}subset \textbackslash{}\textbackslash{}mathbb\textbackslash{}{R\textbackslash{}}\textbackslash{}textasciicircum{}\textbackslash{}{10\textbackslash{}}\textbackslash{}$ be a lattice. Let \textbackslash{}$\textbackslash{}\textbackslash{}mu(\textbackslash{}\textbackslash{}Lambda) = \textbackslash{}\textbackslash{}min\textbackslash{}_\textbackslash{}{\textbackslash{}\textbackslash{}mathbf\textbackslash{}{v\textbackslash{}} \textbackslash{}\textbackslash{}in \textbackslash{}\textbackslash{}Lambda \textbackslash{}\textbackslash{}setminus \textbackslash{}\textbackslash{}\textbackslash{}{0\textbackslash{}\textbackslash{}\textbackslash{}}\textbackslash{}} \textbackslash{}\textbackslash{}|\textbackslash{}\textbackslash{}mathbf\textbackslash{}{v\textbackslash{}}\textbackslash{}\textbackslash{}|\textbackslash{}textasciicircum{}2\textbackslash{}$ be its minimal squared norm and \textbackslash{}$k(\textbackslash{}\textbackslash{}Lambda) = |\textbackslash{}\textbackslash{}\textbackslash{}{\textbackslash{}\textbackslash{}mathbf\textbackslash{}{v\textbackslash{}} \textbackslash{}\textbackslash{}in \textbackslash{}\textbackslash{}Lambda : \textbackslash{}\textbackslash{}|\textbackslash{}\textbackslash{}mathbf\textbackslash{}{v\textbackslash{}}\textbackslash{}\textbackslash{}|\textbackslash{}textasciicircum{}2 = \textbackslash{}\textbackslash{}mu(\textbackslash{}\textbackslash{}Lambda)\textbackslash{}\textbackslash{}\textbackslash{}}|\textbackslash{}$ be its kissing number. If \textbackslash{}$\textbackslash{}\textbackslash{}mu(\textbackslash{}\textbackslash{}Lambda) = 4\textbackslash{}$ and \textbackslash{}$k(\textbackslash{}\textbackslash{}Lambda) = 336\textbackslash{}$, then the determinant of any Gram matrix for \textbackslash{}$\textbackslash{}\textbackslash{}Lambda\textbackslash{}$, 
\end{conjecture}

\begin{lstlisting}[caption={Lean 4 Sketch (v14 Galois)}]
import Mathlib.Analysis.InnerProductSpace.Lattice
import Mathlib.LinearAlgebra.Gram

open Real InnerProductSpace

-- Using Fin 10 for the 10-dimensional Euclidean space
abbrev E10 := EuclideanSpace ℝ (Fin 10)

-- Define minimal squared norm for a lattice Λ
noncomputable def minNormSq (Λ : AddSubgroup E10) : ℝ :=
  inf { ‖v‖^2 | (v : E10) ∈ Λ ∧ v ≠ 0 }

-- Define kissing number for a lattice Λ
-- Assumes the set of minimal vectors is finite, which is true for lattices.
noncomputable def kissingNumber (Λ : AddSubgroup E10) : ℕ :=
  Nat.card { v : Λ | ‖v‖^2 = minNormSq Λ }

-- Define what it mean
\end{lstlisting}

\hrule\medskip

\subsection*{\texttt{periodic\textbackslash{}_packing\textbackslash{}_dim10}}

\textbf{Domain:} discrete\textbackslash{}_geometry \quad \textbf{Class:} 3 \quad \textbf{Verdict:} \textcolor{incompleteamber}{\textbf{INCOMPLETE}}

\textbf{Confidence:} 0.70 \quad \textbf{Domain Prior:} 0.20 \quad \textbf{Sorry gaps:} 4

\textcolor{sorryred}{\textbf{Zero-Sorry Guillotine Applied:} 4 \texttt{sorry} stubs detected. Verdict downgraded from VERIFIED to INCOMPLETE.}

\begin{conjecture}
Let \textbackslash{}$\textbackslash{}\textbackslash{}mathcal\textbackslash{}{P\textbackslash{}}\textbackslash{}_n\textbackslash{}$ be the set of all periodic sphere packings in \textbackslash{}$\textbackslash{}\textbackslash{}mathbb\textbackslash{}{R\textbackslash{}}\textbackslash{}textasciicircum{}n\textbackslash{}$, and let \textbackslash{}$\textbackslash{}\textbackslash{}Delta\textbackslash{}_n = \textbackslash{}\textbackslash{}sup\textbackslash{}_\textbackslash{}{P \textbackslash{}\textbackslash{}in \textbackslash{}\textbackslash{}mathcal\textbackslash{}{P\textbackslash{}}\textbackslash{}_n\textbackslash{}} \textbackslash{}\textbackslash{}text\textbackslash{}{density\textbackslash{}}(P)\textbackslash{}$ be the maximum possible density. Let \textbackslash{}$\textbackslash{}\textbackslash{}Lambda\textbackslash{}_9 \textbackslash{}\textbackslash{}subset \textbackslash{}\textbackslash{}mathbb\textbackslash{}{R\textbackslash{}}\textbackslash{}textasciicircum{}9\textbackslash{}$ be the optimal lattice for sphere packing in 9 dimensions (conjecturally, the Coxeter-Todd lattice). We conjecture that \textbackslash{}$\textbackslash{}\textbackslash{}Delta\textbackslash{}_\textbackslash{}{10\textbackslash{}}\textbackslash{}$ is achieved not by a unique packing, 
\end{conjecture}

\begin{lstlisting}[caption={Lean 4 Sketch (v14 Galois)}]
import Mathlib.Geometry.Euclidean.SpherePacking.Basic
import Mathlib.Analysis.InnerProductSpace.PiL2
import Mathlib.LinearAlgebra.Matrix.Gram

open Set Metric

/- We model R^n as `EuclideanSpace ℝ (Fin n)` -/

/-- A predicate asserting that a sphere packing `P` in `n+1` dimensions is a
    periodic stacking of a layer packing `P_layer` in `n` dimensions. -/
def IsPeriodicStacking {n : ℕ} (P : SpherePacking (Fin (n+1))) (P_layer : SpherePacking (Fin n)) : Prop := sorry
  -- sorry justification: This definition is non-trivial. It requires formalizing the geometric
  -- concept of slicing a point
\end{lstlisting}

\hrule\medskip

\subsection*{\texttt{bessel\textbackslash{}_moment\textbackslash{}_c6\textbackslash{}_0}}

\textbf{Domain:} special\textbackslash{}_functions \quad \textbf{Class:} 3 \quad \textbf{Verdict:} \textcolor{incompleteamber}{\textbf{INCOMPLETE}}

\textbf{Confidence:} 0.70 \quad \textbf{Domain Prior:} 0.25 \quad \textbf{Sorry gaps:} 6

\textcolor{sorryred}{\textbf{Zero-Sorry Guillotine Applied:} 6 \texttt{sorry} stubs detected. Verdict downgraded from VERIFIED to INCOMPLETE.}

\begin{conjecture}
Let \textbackslash{}$L\textbackslash{}_n := \textbackslash{}\textbackslash{}int\textbackslash{}_0\textbackslash{}textasciicircum{}\textbackslash{}\textbackslash{}infty x K\textbackslash{}_0(x)\textbackslash{}textasciicircum{}n dx\textbackslash{}$ for \textbackslash{}$n \textbackslash{}\textbackslash{}in \textbackslash{}\textbackslash{}mathbb\textbackslash{}{N\textbackslash{}}\textbackslash{}$, where \textbackslash{}$K\textbackslash{}_0\textbackslash{}$ is the modified Bessel function of the second kind of order zero. Then the following equality holds:\textbackslash{}n\textbackslash{}$\textbackslash{}$ L\textbackslash{}_6 = \textbackslash{}\textbackslash{}frac\textbackslash{}{3645\textbackslash{}}\textbackslash{}{128\textbackslash{}\textbackslash{}pi\textbackslash{}textasciicircum{}2\textbackslash{}} \textbackslash{}\textbackslash{}zeta(3) L(\textbackslash{}\textbackslash{}chi\textbackslash{}_\textbackslash{}{-3\textbackslash{}}, 4) \textbackslash{}$\textbackslash{}$ \textbackslash{}nwhere \textbackslash{}$\textbackslash{}\textbackslash{}zeta\textbackslash{}$ is the Riemann zeta function and \textbackslash{}$L(\textbackslash{}\textbackslash{}chi\textbackslash{}_\textbackslash{}{-3\textbackslash{}}, s) := \textbackslash{}\textbackslash{}sum\textbackslash{}_\textbackslash{}{k=1\textbackslash{}}\textbackslash{}textasciicircum{}\textbackslash{}\textbackslash{}infty \textbackslash{}\textbackslash{}frac\textbackslash{}{\textbackslash{}\textbackslash{}chi\textbackslash{}_\textbackslash{}{-3\textbackslash{}}(k)\textbackslash{}}\textbackslash{}{k\textbackslash{}textasciicircum{}s\textbackslash{}}\textbackslash{}$ is the Dirichlet L-function for the
\end{conjecture}

\begin{lstlisting}[caption={Lean 4 Sketch (v14 Galois)}]
import Analysis.SpecialFunctions.Bessel
import Analysis.SpecialFunctions.Zeta
import Analysis.SpecialFunctions.Trigonometric.Core
import NumberTheory.DirichletCharacter
import Mathlib.Data.Real.Pi

open Real Nat Filter Asymptotics

/- Define the Dirichlet L-function for the primitive character mod 3, χ₃(k) = (k/3) -/
noncomputable def dirichletL_chi_neg_3 (s : ℝ) : ℝ := ∑' k, (if k % 3 = 1 then 1 else if k % 3 = 2 then -1 else 0) / (k : ℝ) ^ s

-- Conjectured property: Convergence for s > 0
-- theorem dirichletL_chi_neg_3_converges {s : ℝ} (hs : 0 < s) : Summable (fun k ↦ (if k % 3 = 1 then 1 
\end{lstlisting}

\hrule\medskip

\subsection*{\texttt{feynman\textbackslash{}_3loop\textbackslash{}_sunrise}}

\textbf{Domain:} mathematical\textbackslash{}_physics \quad \textbf{Class:} 4 \quad \textbf{Verdict:} \textcolor{refutedred}{\textbf{REFUTED}}

\textbf{Confidence:} 0.57 \quad \textbf{Domain Prior:} 0.15 \quad \textbf{Sorry gaps:} 8

\begin{conjecture}
Let \textbackslash{}$\textbackslash{}\textbackslash{}\textbackslash{}\textbackslash{}mathcal\textbackslash{}{J\textbackslash{}}\textbackslash{}textasciicircum{}\textbackslash{}{(0)\textbackslash{}}(t)\textbackslash{}$ be the renormalized finite part of the dimensionless equal-mass three-loop sunrise integral, where \textbackslash{}$t=p\textbackslash{}textasciicircum{}2/m\textbackslash{}textasciicircum{}2\textbackslash{}$. Let \textbackslash{}$f \textbackslash{}\textbackslash{}\textbackslash{}\textbackslash{}in S\textbackslash{}_4(\textbackslash{}\textbackslash{}\textbackslash{}\textbackslash{}Gamma\textbackslash{}_0(6))\textbackslash{}$ be the unique normalized newform with Fourier expansion \textbackslash{}$f(z) = q - 4q\textbackslash{}textasciicircum{}2 - 3q\textbackslash{}textasciicircum{}3 + O(q\textbackslash{}textasciicircum{}4)\textbackslash{}$, and let \textbackslash{}$L(f,s)\textbackslash{}$ be its associated L-series. There exists a rational number \textbackslash{}$r \textbackslash{}\textbackslash{}in \textbackslash{}\textbackslash{}mathbb\textbackslash{}{Q\textbackslash{}}\textbackslash{}$ such that the value of the integral
\end{conjecture}

\begin{lstlisting}[caption={Lean 4 Sketch (v14 Galois)}]
import Analysis.SpecialFunctions.LSeries\
import NumberTheory.ModularForms\
import Analysis.SpecialFunctions.Zeta\
import NumberTheory.Clausen\
\
open Complex Topos\
\
-- We postulate the existence of the finite part of the three-loop sunrise integral.\
-- A full definition would require the theory of dimensional regularization and renormalization.\
noncomputable def threeLoopSunriseFinitePart (t : ℝ) : ℝ := sorry\
\
-- Let f be the unique newform in S_4(Γ₀(6)). Mathlib's modular form library\
-- can in principle define this object.\
def Gamma0_6_weight4_newform : CuspForm (Gamma0 6) 4 := sorr
\end{lstlisting}

\hrule\medskip

\subsection*{\texttt{townes\textbackslash{}_soliton}}

\textbf{Domain:} mathematical\textbackslash{}_physics \quad \textbf{Class:} 4 \quad \textbf{Verdict:} \textcolor{incompleteamber}{\textbf{INCOMPLETE}}

\textbf{Confidence:} 0.70 \quad \textbf{Domain Prior:} 0.15 \quad \textbf{Sorry gaps:} 2

\textcolor{sorryred}{\textbf{Zero-Sorry Guillotine Applied:} 2 \texttt{sorry} stubs detected. Verdict downgraded from VERIFIED to INCOMPLETE.}

\begin{conjecture}
Let \textbackslash{}$G=(V,E)\textbackslash{}$ be a finite simple connected graph with \textbackslash{}$n=|V|\textbackslash{}$ vertices and \textbackslash{}$m=|E|\textbackslash{}$ edges. Let \textbackslash{}$L\textbackslash{}$ be the normalized graph Laplacian of \textbackslash{}$G\textbackslash{}$. For a given mass \textbackslash{}$M>0\textbackslash{}$, consider the energy functional \textbackslash{}$E: \textbackslash{}\textbackslash{}mathbb\textbackslash{}{R\textbackslash{}}\textbackslash{}textasciicircum{}n \textbackslash{}\textbackslash{}to \textbackslash{}\textbackslash{}mathbb\textbackslash{}{R\textbackslash{}}\textbackslash{}$ defined by \textbackslash{}$E(u) = \textbackslash{}\textbackslash{}frac\textbackslash{}{1\textbackslash{}}\textbackslash{}{2\textbackslash{}} \textbackslash{}\textbackslash{}sum\textbackslash{}_\textbackslash{}{(i,j) \textbackslash{}\textbackslash{}in E\textbackslash{}} (u\textbackslash{}_i - u\textbackslash{}_j)\textbackslash{}textasciicircum{}2 - \textbackslash{}\textbackslash{}frac\textbackslash{}{1\textbackslash{}}\textbackslash{}{4\textbackslash{}} \textbackslash{}\textbackslash{}sum\textbackslash{}_\textbackslash{}{v \textbackslash{}\textbackslash{}in V\textbackslash{}} u\textbackslash{}_v\textbackslash{}textasciicircum{}4\textbackslash{}$, restricted to the sphere \textbackslash{}$S\textbackslash{}_M = \textbackslash{}\textbackslash{}\textbackslash{}{ u \textbackslash{}\textbackslash{}in \textbackslash{}\textbackslash{}mathbb\textbackslash{}{R\textbackslash{}}\textbackslash{}textasciicircum{}n : \textbackslash{}\textbackslash{}|
\end{conjecture}

\begin{lstlisting}[caption={Lean 4 Sketch (v14 Galois)}]
import Mathlib.Combinatorics.SimpleGraph.Connectivity
import Mathlib.Analysis.Calculus.FDeriv.Manifold
import Mathlib.Geometry.Manifold.Instances.Sphere
import Mathlib.Topology.Homology.Simplicial

/- We model the space of functions on the vertices of G as a finite-dimensional real vector space. -/
variable {V : Type*} [Fintype V] [DecidableEq V]
variable (G : SimpleGraph V)

/- The energy functional for the discrete nonlinear Schrödinger equation (DNLS). -/
noncomputable def dnlsEnergy (u : V → ℝ) : ℝ :=
  (1/2) * (∑ e in G.edgeSet, (u e.val.1 - u e.val.2)^2) - (1/4) * (∑ v, (u v)^4)

/- The 
\end{lstlisting}

\hrule\medskip

\subsection*{\texttt{elliptic\textbackslash{}_curve\textbackslash{}_rank\textbackslash{}_30}}

\textbf{Domain:} number\textbackslash{}_theory \quad \textbf{Class:} 4 \quad \textbf{Verdict:} \textcolor{incompleteamber}{\textbf{INCOMPLETE}}

\textbf{Confidence:} 0.70 \quad \textbf{Domain Prior:} 0.05 \quad \textbf{Sorry gaps:} 1

\textcolor{sorryred}{\textbf{Zero-Sorry Guillotine Applied:} 1 \texttt{sorry} stubs detected. Verdict downgraded from VERIFIED to INCOMPLETE.}

\begin{conjecture}
Let \textbackslash{}$E\textbackslash{}_d\textbackslash{}$ be the family of elliptic curves over \textbackslash{}$\textbackslash{}\textbackslash{}mathbb\textbackslash{}{Q\textbackslash{}}\textbackslash{}$ defined by the equation \textbackslash{}$y\textbackslash{}textasciicircum{}2 = x\textbackslash{}textasciicircum{}3 - dx\textbackslash{}$. Let \textbackslash{}$P\textbackslash{}_4 = 2 \textbackslash{}\textbackslash{}cdot 3 \textbackslash{}\textbackslash{}cdot 5 \textbackslash{}\textbackslash{}cdot 7 = 210\textbackslash{}$ be the fourth primorial. The sum of the Mordell-Weil ranks of the twists of \textbackslash{}$E\textbackslash{}_1\textbackslash{}$ by the positive square-free divisors of \textbackslash{}$P\textbackslash{}_4\textbackslash{}$ is exactly 30. Formally, let \textbackslash{}$S = \textbackslash{}\textbackslash{}\textbackslash{}{d \textbackslash{}\textbackslash{}in \textbackslash{}\textbackslash{}mathbb\textbackslash{}{N\textbackslash{}}\textbackslash{}textasciicircum{}* \textbackslash{}\textbackslash{}mid d \textbackslash{}\textbackslash{}text\textbackslash{}{ is square-free and \textbackslash{}} d|P\textbackslash{}_4\textbackslash{}\textbackslash{}\textbackslash{}}\textbackslash{}$. Then: \textbackslash{}$\textbackslash{}$\textbackslash{}\textbackslash{}sum\textbackslash{}_\textbackslash{}{
\end{conjecture}

\begin{lstlisting}[caption={Lean 4 Sketch (v14 Galois)}]
import Mathlib.NumberTheory.EllipticCurve.Rank
import Mathlib.Data.Nat.Squarefree

open EllipticCurve Finset

-- Define the family of elliptic curves E_d: y^2 = x^3 - d*x
def E_twist (d : ℤ) : EllipticCurve ℚ :=
  EllipticCurve.of_ainvs 0 (-(d : ℚ)) 0 0 0

-- The conjecture states that the sum of ranks of twists of E_1 by the
-- positive square-free divisors of the 4th primorial (210) is 30.
-- Since 210 is square-free, its divisors are all square-free.
theorem galois_conjecture_elliptic_curve_rank_30 : 
  ∑ d in Nat.divisors 210, rank (E_twist d) = 30 := by
  -- This identity is computational
\end{lstlisting}

\hrule\medskip

\subsection*{\texttt{elliptic\textbackslash{}_curve\textbackslash{}_rank\textbackslash{}_torsion\textbackslash{}_z7z}}

\textbf{Domain:} number\textbackslash{}_theory \quad \textbf{Class:} 4 \quad \textbf{Verdict:} \textcolor{incompleteamber}{\textbf{INCOMPLETE}}

\textbf{Confidence:} 0.70 \quad \textbf{Domain Prior:} 0.05 \quad \textbf{Sorry gaps:} 2

\textcolor{sorryred}{\textbf{Zero-Sorry Guillotine Applied:} 2 \texttt{sorry} stubs detected. Verdict downgraded from VERIFIED to INCOMPLETE.}

\begin{conjecture}
Let \textbackslash{}$K\textbackslash{}$ be a number field and \textbackslash{}$E\textbackslash{}$ be an elliptic curve defined over \textbackslash{}$K\textbackslash{}$. If the torsion subgroup of the Mordell-Weil group \textbackslash{}$E(K)\textbackslash{}$ contains a subgroup isomorphic to \textbackslash{}$\textbackslash{}\textbackslash{}mathbb\textbackslash{}{Z\textbackslash{}}/7\textbackslash{}\textbackslash{}mathbb\textbackslash{}{Z\textbackslash{}} \textbackslash{}\textbackslash{}times \textbackslash{}\textbackslash{}mathbb\textbackslash{}{Z\textbackslash{}}/7\textbackslash{}\textbackslash{}mathbb\textbackslash{}{Z\textbackslash{}}\textbackslash{}$, then the rank of \textbackslash{}$E(K)\textbackslash{}$ is zero. Formally, if \textbackslash{}$(\textbackslash{}\textbackslash{}mathbb\textbackslash{}{Z\textbackslash{}}/7\textbackslash{}\textbackslash{}mathbb\textbackslash{}{Z\textbackslash{}})\textbackslash{}textasciicircum{}2 \textbackslash{}\textbackslash{}hookrightarrow E(K)\textbackslash{}_\textbackslash{}{\textbackslash{}\textbackslash{}text\textbackslash{}{tors\textbackslash{}}\textbackslash{}}\textbackslash{}$, then \textbackslash{}$\textbackslash{}\textbackslash{}text\textbackslash{}{rank\textbackslash{}}\textbackslash{}_\textbackslash{}{\textbackslash{}\textbackslash{}mathbb\textbackslash{}{Z\textbackslash{}}\textbackslash{}}(E(K)) = 0\textbackslash{}$.
\end{conjecture}

\begin{lstlisting}[caption={Lean 4 Sketch (v14 Galois)}]
import Mathlib.NumberTheory.EllipticCurve.Point
import Mathlib.NumberTheory.NumberField.Basic
import Mathlib.GroupTheory.Finiteness
import Mathlib.Algebra.Module.Rank
import Mathlib.GroupTheory.SpecificGroups.Cyclic

open EllipticCurve Additive

-- This conjecture proposes a strong structural property for a specific class of elliptic curves.
-- It connects the torsion structure to the rank of the Mordell-Weil group.
-- A proof would be a major result in the arithmetic of elliptic curves.
-- We use `Module.rank ℤ (Point E)` for the Mordell-Weil rank of E(K).

-- We define a placeholder for the 
\end{lstlisting}

\hrule\medskip

\subsection*{\texttt{mzv\textbackslash{}_decomposition\textbackslash{}_c5}}

\textbf{Domain:} number\textbackslash{}_theory \quad \textbf{Class:} 4 \quad \textbf{Verdict:} \textcolor{incompleteamber}{\textbf{INCOMPLETE}}

\textbf{Confidence:} 0.70 \quad \textbf{Domain Prior:} 0.05 \quad \textbf{Sorry gaps:} 6

\textcolor{sorryred}{\textbf{Zero-Sorry Guillotine Applied:} 6 \texttt{sorry} stubs detected. Verdict downgraded from VERIFIED to INCOMPLETE.}

\begin{conjecture}
Let \textbackslash{}$I\textbackslash{}_k\textbackslash{}textasciicircum{}d = \textbackslash{}\textbackslash{}\textbackslash{}{(k\textbackslash{}_1, \textbackslash{}\textbackslash{}dots, k\textbackslash{}_d) \textbackslash{}\textbackslash{}in \textbackslash{}\textbackslash{}mathbb\textbackslash{}{Z\textbackslash{}}\textbackslash{}textasciicircum{}d \textbackslash{}\textbackslash{}mid \textbackslash{}\textbackslash{}sum\textbackslash{}_\textbackslash{}{i=1\textbackslash{}}\textbackslash{}textasciicircum{}d k\textbackslash{}_i = k, k\textbackslash{}_1 > 1, k\textbackslash{}_i \textbackslash{}\textbackslash{}ge 1 \textbackslash{}\textbackslash{}\textbackslash{}}\textbackslash{}$. The sum of all admissible multiple zeta values of weight 5 with depth greater than or equal to 2 is equal to \textbackslash{}$2\textbackslash{}\textbackslash{}zeta(5)\textbackslash{}$. Formally, \textbackslash{}$\textbackslash{}$ \textbackslash{}\textbackslash{}sum\textbackslash{}_\textbackslash{}{d=2\textbackslash{}}\textbackslash{}textasciicircum{}\textbackslash{}{4\textbackslash{}} \textbackslash{}\textbackslash{}sum\textbackslash{}_\textbackslash{}{\textbackslash{}\textbackslash{}mathbf\textbackslash{}{k\textbackslash{}} \textbackslash{}\textbackslash{}in I\textbackslash{}_5\textbackslash{}textasciicircum{}d\textbackslash{}} \textbackslash{}\textbackslash{}zeta(\textbackslash{}\textbackslash{}mathbf\textbackslash{}{k\textbackslash{}}) = 2\textbackslash{}\textbackslash{}zeta(5) \textbackslash{}$\textbackslash{}$
\end{conjecture}

\begin{lstlisting}[caption={Lean 4 Sketch (v14 Galois)}]
import Mathlib.Analysis.SpecialFunctions.Zeta
import Mathlib.Data.Real.Basic

-- We postulate the existence of a Multiple Zeta Value function and its fundamental properties.
-- A full formalization would require a significant library development.

-- Represents ζ(k₁, ..., kᵣ) for k = [k₁, ..., kᵣ]
opaque mzv (k : List ℕ) : ℝ

-- Postulate: The sum of depth-2 MZVs of weight 5 is ζ(5).
-- This is a specific instance of the Sum Formula.
theorem mzv_sum_formula_d2_w5 : mzv [4,1] + mzv [3,2] + mzv [2,3] = Real.zeta 5 := sorry

-- Postulate: Duality relations for weight 5.
-- These are instances of 
\end{lstlisting}

\hrule\medskip

\subsection*{\texttt{tracy\textbackslash{}_widom\textbackslash{}_f2\textbackslash{}_mean}}

\textbf{Domain:} mathematical\textbackslash{}_physics \quad \textbf{Class:} 3 \quad \textbf{Verdict:} \textcolor{incompleteamber}{\textbf{INCOMPLETE}}

\textbf{Confidence:} 0.70 \quad \textbf{Domain Prior:} 0.15 \quad \textbf{Sorry gaps:} 2

\textcolor{sorryred}{\textbf{Zero-Sorry Guillotine Applied:} 2 \texttt{sorry} stubs detected. Verdict downgraded from VERIFIED to INCOMPLETE.}

\begin{conjecture}
Let \textbackslash{}$F\textbackslash{}_2(s)\textbackslash{}$ be the Tracy-Widom GUE distribution function. The mean of this distribution, \textbackslash{}$\textbackslash{}\textbackslash{}mu\textbackslash{}_2 = \textbackslash{}\textbackslash{}int\textbackslash{}_\textbackslash{}{-\textbackslash{}\textbackslash{}infty\textbackslash{}}\textbackslash{}textasciicircum{}\textbackslash{}{\textbackslash{}\textbackslash{}infty\textbackslash{}} s \textbackslash{}\textbackslash{}, dF\textbackslash{}_2(s)\textbackslash{}$, is given by the exact expression: \textbackslash{}\textbackslash{}n\textbackslash{}\textbackslash{}[ \textbackslash{}\textbackslash{}mu\textbackslash{}_2 = - \textbackslash{}\textbackslash{}left( \textbackslash{}\textbackslash{}frac\textbackslash{}{3\textbackslash{}\textbackslash{}pi \textbackslash{}\textbackslash{}zeta(3)\textbackslash{}}\textbackslash{}{2\textbackslash{}} \textbackslash{}\textbackslash{}right)\textbackslash{}textasciicircum{}\textbackslash{}{1/3\textbackslash{}} \textbackslash{}\textbackslash{}]\textbackslash{}nwhere \textbackslash{}$\textbackslash{}\textbackslash{}zeta(s)\textbackslash{}$ is the Riemann zeta function.
\end{conjecture}

\begin{lstlisting}[caption={Lean 4 Sketch (v14 Galois)}]
import Mathlib.Analysis.SpecialFunctions.Zeta
import Mathlib.Analysis.SpecialFunctions.Airy
import Mathlib.Analysis.Calculus.ParametricIntegral

open Real Set Filter Topology

-- We postulate the existence and properties of the Hastings-McLeod solution
-- to the Painlevé II equation q'' = 2q^3 + xq.
postulate q_HM (x : ℝ) : ℝ
postulate q_HM_is_solution : ∀ x, Differentiable.DifferentiableAt ℝ (fun y ↦ deriv q_HM y) x ∧
  deriv (deriv q_HM) x = 2 * (q_HM x)^3 + x * q_HM x
postulate q_HM_asymptotics : (fun x ↦ q_HM x / airyAi x) →* (nhds_atTop) (const ℝ 1)

-- We also need integrability properti
\end{lstlisting}

\hrule\medskip

\subsection*{\texttt{tracy\textbackslash{}_widom\textbackslash{}_f2\textbackslash{}_variance}}

\textbf{Domain:} mathematical\textbackslash{}_physics \quad \textbf{Class:} 3 \quad \textbf{Verdict:} \textcolor{incompleteamber}{\textbf{INCOMPLETE}}

\textbf{Confidence:} 0.70 \quad \textbf{Domain Prior:} 0.15 \quad \textbf{Sorry gaps:} 3

\textcolor{sorryred}{\textbf{Zero-Sorry Guillotine Applied:} 3 \texttt{sorry} stubs detected. Verdict downgraded from VERIFIED to INCOMPLETE.}

\begin{conjecture}
Let \textbackslash{}$\textbackslash{}\textbackslash{}mu\textbackslash{}_1, \textbackslash{}\textbackslash{}sigma\textbackslash{}_1\textbackslash{}textasciicircum{}2\textbackslash{}$ be the mean and variance of the GUE Tracy-Widom distribution \textbackslash{}$F\textbackslash{}_1\textbackslash{}$, and let \textbackslash{}$\textbackslash{}\textbackslash{}mu\textbackslash{}_2, \textbackslash{}\textbackslash{}sigma\textbackslash{}_2\textbackslash{}textasciicircum{}2\textbackslash{}$ be the mean and variance of the GOE Tracy-Widom distribution \textbackslash{}$F\textbackslash{}_2\textbackslash{}$. Then these quantities are related by the linear identity:\textbackslash{}\textbackslash{}n\textbackslash{}\textbackslash{}[ \textbackslash{}\textbackslash{}sigma\textbackslash{}_2\textbackslash{}textasciicircum{}2 + 2\textbackslash{}\textbackslash{}mu\textbackslash{}_2 + \textbackslash{}\textbackslash{}sigma\textbackslash{}_1\textbackslash{}textasciicircum{}2 = 0 \textbackslash{}\textbackslash{}]
\end{conjecture}

\begin{lstlisting}[caption={Lean 4 Sketch (v14 Galois)}]
import Mathlib.Analysis.SpecialFunctions.Integrals
import Mathlib.Probability.Theory.Variance

-- We axiomatically define the Tracy-Widom distributions as they are not yet in Mathlib.
-- A full definition would require the theory of Fredholm determinants of Airy kernel operators.
axiom tracyWidomF1 : Measure ℝ
axiom tracyWidomF2 : Measure ℝ

-- We assume they are probability measures with finite second moments.
@[instance] axiom F1_isProb : IsProbabilityMeasure tracyWidomF1
@[instance] axiom F2_isProb : IsProbabilityMeasure tracyWidomF2
axiom F1_finite_moment_2 : (∫⁻ x, ‖x^2‖₊ ∂tracyWidomF1) <
\end{lstlisting}

\hrule\medskip

\subsection*{\texttt{crossing\textbackslash{}_number\textbackslash{}_kn}}

\textbf{Domain:} combinatorics \quad \textbf{Class:} 4 \quad \textbf{Verdict:} \textcolor{incompleteamber}{\textbf{INCOMPLETE}}

\textbf{Confidence:} 0.70 \quad \textbf{Domain Prior:} 0.35 \quad \textbf{Sorry gaps:} 2

\textcolor{sorryred}{\textbf{Zero-Sorry Guillotine Applied:} 2 \texttt{sorry} stubs detected. Verdict downgraded from VERIFIED to INCOMPLETE.}

\begin{conjecture}
Let \textbackslash{}$\textbackslash{}\textbackslash{}text\textbackslash{}{cr\textbackslash{}}(K\textbackslash{}_n)\textbackslash{}$ denote the crossing number of the complete graph on \textbackslash{}$n\textbackslash{}$ vertices. For any integer \textbackslash{}$n \textbackslash{}\textbackslash{}ge 3\textbackslash{}$, the following inequality holds: \textbackslash{}\textbackslash{}\textbackslash{}\textbackslash{}\textbackslash{}$\textbackslash{}$ \textbackslash{}\textbackslash{}text\textbackslash{}{cr\textbackslash{}}(K\textbackslash{}_\textbackslash{}{n+1\textbackslash{}}) \textbackslash{}\textbackslash{}ge \textbackslash{}\textbackslash{}text\textbackslash{}{cr\textbackslash{}}(K\textbackslash{}_n) + \textbackslash{}{\textbackslash{}\textbackslash{}lfloor n/2 \textbackslash{}\textbackslash{}rfloor \textbackslash{}\textbackslash{}choose 2\textbackslash{}} \textbackslash{}\textbackslash{}left\textbackslash{}\textbackslash{}lfloor \textbackslash{}\textbackslash{}frac\textbackslash{}{n-1\textbackslash{}}\textbackslash{}{2\textbackslash{}} \textbackslash{}\textbackslash{}right\textbackslash{}\textbackslash{}rfloor. \textbackslash{}\textbackslash{}\textbackslash{}$\textbackslash{}$
\end{conjecture}

\begin{lstlisting}[caption={Lean 4 Sketch (v14 Galois)}]
import Mathlib.Data.Nat.Choose.Basic
import Mathlib.Tactic

-- We postulate the existence of the crossing number function for the complete graph K_n.
-- In a full formalization, this would be defined as the minimum number of crossings
-- over all planar drawings of K_n.
opaque cr_K (n : ℕ) : ℕ

-- Postulate a basic, known property that cr_K is non-decreasing.
axiom cr_K_mono : Monotone cr_K

/-- The conjectured incremental lower bound for `cr(K_{n+1}) - cr(K_n)`.
This value is equivalent to `Z(n+1) - Z(n)`, where `Z` is the conjectured
formula for `cr(K_n)`. -/
def incremental_crossing_lower_b
\end{lstlisting}

\hrule\medskip

\subsection*{\texttt{kcore\textbackslash{}_threshold\textbackslash{}_c3}}

\textbf{Domain:} combinatorics \quad \textbf{Class:} 3 \quad \textbf{Verdict:} \textcolor{incompleteamber}{\textbf{INCOMPLETE}}

\textbf{Confidence:} 0.70 \quad \textbf{Domain Prior:} 0.35 \quad \textbf{Sorry gaps:} 4

\textcolor{sorryred}{\textbf{Zero-Sorry Guillotine Applied:} 4 \texttt{sorry} stubs detected. Verdict downgraded from VERIFIED to INCOMPLETE.}

\begin{conjecture}
Let \textbackslash{}$G(n, p)\textbackslash{}$ be the Erdős–Rényi random graph with \textbackslash{}$p=c/n\textbackslash{}$. Let \textbackslash{}$C\textbackslash{}_k(G)\textbackslash{}$ be the \textbackslash{}$k\textbackslash{}$-core of \textbackslash{}$G\textbackslash{}$. Let \textbackslash{}$C\textbackslash{}_\textbackslash{}\textbackslash{}triangle(G)\textbackslash{}$ be the triangle-core of \textbackslash{}$G\textbackslash{}$, defined as the largest induced subgraph of \textbackslash{}$G\textbackslash{}$ where every vertex is part of at least one triangle. Let \textbackslash{}$c\textbackslash{}_k\textbackslash{}$ and \textbackslash{}$c\textbackslash{}_\textbackslash{}\textbackslash{}triangle\textbackslash{}$ be the respective threshold constants for the emergence of a non-empty \textbackslash{}$C\textbackslash{}_k(G)\textbackslash{}$ and \textbackslash{}$C\textbackslash{}_\textbackslash{}\textbackslash{}triangle(G)\textbackslash{}$. The conjecture is th
\end{conjecture}

\begin{lstlisting}[caption={Lean 4 Sketch (v14 Galois)}]
import analysis.special_functions.exp_log
import combinatorics.simple_graph.core
import probability_theory.borel_cantelli

noncomputable theory
open filter

/- We assume a framework for G(n,p) random graphs, providing a probability space
   on `simple_graph (fin n)` for each n. -/
axiom gnp (n : ℕ) (p : ℝ) : Type*
axiom gnp_prob {n : ℕ} {p : ℝ} : measure_space (gnp n p)

-- Definition of the triangle-core vertices via a peeling operator
def is_in_triangle {V : Type*} (G : simple_graph V) (v : V) (S : set V) : Prop :=
  ∃ u w ∈ S, G.adj v u ∧ G.adj v w ∧ G.adj u w

def triangle_core_peeling_ope
\end{lstlisting}

\hrule\medskip

\subsection*{\texttt{covering\textbackslash{}_C13\textbackslash{}_k7\textbackslash{}_t4}}

\textbf{Domain:} combinatorics \quad \textbf{Class:} 3 \quad \textbf{Verdict:} \textcolor{incompleteamber}{\textbf{INCOMPLETE}}

\textbf{Confidence:} 0.70 \quad \textbf{Domain Prior:} 0.35 \quad \textbf{Sorry gaps:} 7

\textcolor{sorryred}{\textbf{Zero-Sorry Guillotine Applied:} 7 \texttt{sorry} stubs detected. Verdict downgraded from VERIFIED to INCOMPLETE.}

\begin{conjecture}
Let \textbackslash{}$V\textbackslash{}$ be a set of 13 points, and \textbackslash{}$\textbackslash{}\textbackslash{}\textbackslash{}\textbackslash{}mathcal\textbackslash{}{B\textbackslash{}}\textbackslash{}$ be a collection of 7-subsets of \textbackslash{}$V\textbackslash{}$ (called blocks) such that every 4-subset of \textbackslash{}$V\textbackslash{}$ is contained in at least one block of \textbackslash{}$\textbackslash{}\textbackslash{}\textbackslash{}\textbackslash{}mathcal\textbackslash{}{B\textbackslash{}}\textbackslash{}$. Let \textbackslash{}$C(13, 7, 4)\textbackslash{}$ be the minimum possible size of such a collection \textbackslash{}$\textbackslash{}\textbackslash{}\textbackslash{}\textbackslash{}mathcal\textbackslash{}{B\textbackslash{}}\textbackslash{}$. The conjecture is twofold: (1) \textbackslash{}$C(13, 7, 4) = 23\textbackslash{}$. (2) For any such minimal collection \textbackslash{}$\textbackslash{}\textbackslash{}\textbackslash{}\textbackslash{}mathcal\textbackslash{}{B\textbackslash{}}\textbackslash{}$ of size 23, the mu
\end{conjecture}

\begin{lstlisting}[caption={Lean 4 Sketch (v14 Galois)}]
import Mathlib.Data.Finset.Card\
import Mathlib.Data.Fintype.Card\
import Mathlib.Data.Multiset.Basic\
\
open Finset\
\
/-! # Formalization of Covering Designs -/\
\
structure CoveringDesign (n k t : ℕ) where\
  V : Type*\
  [fintypeV : Fintype V]\
  card_V : Fintype.card V = n\
  blocks : Finset (Finset V)\
  card_block_prop : ∀ B ∈ blocks, B.card = k\
  covering_prop : ∀ T : Finset V, T.card = t → ∃ B ∈ blocks, T ⊆ B\
\
def covering_number (n k t : ℕ) : ℕ :=\
  sInf { b : ℕ | ∃ (D : CoveringDesign n k t), D.blocks.card = b }\
\
def replication_number {n k t : ℕ} (D : CoveringDesign n k t) (x
\end{lstlisting}

\hrule\medskip

\subsection*{\texttt{cwcode\textbackslash{}_29\textbackslash{}_8\textbackslash{}_5}}

\textbf{Domain:} coding\textbackslash{}_theory \quad \textbf{Class:} 3 \quad \textbf{Verdict:} \textcolor{incompleteamber}{\textbf{INCOMPLETE}}

\textbf{Confidence:} 0.70 \quad \textbf{Domain Prior:} 0.30 \quad \textbf{Sorry gaps:} 10

\textcolor{sorryred}{\textbf{Zero-Sorry Guillotine Applied:} 10 \texttt{sorry} stubs detected. Verdict downgraded from VERIFIED to INCOMPLETE.}

\begin{conjecture}
Let \textbackslash{}$C\textbackslash{}$ be a binary linear code with parameters \textbackslash{}$[n, k, d] = [29, 8, 5]\textbackslash{}$. Let \textbackslash{}$C\textbackslash{}textasciicircum{}\textbackslash{}\textbackslash{}\textbackslash{}\textbackslash{}perp\textbackslash{}$ be the dual code of \textbackslash{}$C\textbackslash{}$, with minimum distance \textbackslash{}$d\textbackslash{}textasciicircum{}\textbackslash{}\textbackslash{}\textbackslash{}\textbackslash{}perp\textbackslash{}$. Then, \textbackslash{}$d\textbackslash{}textasciicircum{}\textbackslash{}\textbackslash{}\textbackslash{}\textbackslash{}perp = 2\textbackslash{}$.
\end{conjecture}

\begin{lstlisting}[caption={Lean 4 Sketch (v14 Galois)}]
import Mathlib.LinearAlgebra.FiniteDimensional\
import Mathlib.Data.ZMod.Basic\
import Mathlib.Data.Fintype.Card\
import Mathlib.Data.Fin.VecNotation\
\
universe u\
\
-- Define a binary linear code as a subspace of (Fin n → ZMod 2)\
structure LinearCode (n : ℕ) where\
  carrier : Subspace (ZMod 2) (Fin n → ZMod 2)\
\
-- Hamming weight of a vector\
def hammingWeight {n : ℕ} (v : Fin n → ZMod 2) : ℕ :=\
  Fintype.card {i | v i ≠ 0}\
\
-- Minimum distance of a code (excluding the zero vector)\
def minDist (C : LinearCode n) : ℕ :=\
  if h : C.carrier = ⊥ then 0\
  else sInf {w | ∃ v ∈ C.carrier, 
\end{lstlisting}

\hrule\medskip

\subsection*{\texttt{inverse\textbackslash{}_galois\textbackslash{}_m23}}

\textbf{Domain:} number\textbackslash{}_theory \quad \textbf{Class:} 4 \quad \textbf{Verdict:} \textcolor{incompleteamber}{\textbf{INCOMPLETE}}

\textbf{Confidence:} 0.70 \quad \textbf{Domain Prior:} 0.05 \quad \textbf{Sorry gaps:} 1

\textcolor{sorryred}{\textbf{Zero-Sorry Guillotine Applied:} 1 \texttt{sorry} stubs detected. Verdict downgraded from VERIFIED to INCOMPLETE.}

\begin{conjecture}
Let \textbackslash{}$M\textbackslash{}_\textbackslash{}{23\textbackslash{}}\textbackslash{}$ be the Mathieu group on 23 points. There exists an irreducible polynomial \textbackslash{}$f \textbackslash{}\textbackslash{}in \textbackslash{}\textbackslash{}mathbb\textbackslash{}{Q\textbackslash{}}[X]\textbackslash{}$ of degree 23 such that the Galois group of the splitting field of \textbackslash{}$f\textbackslash{}$ over \textbackslash{}$\textbackslash{}\textbackslash{}mathbb\textbackslash{}{Q\textbackslash{}}\textbackslash{}$ is isomorphic to \textbackslash{}$M\textbackslash{}_\textbackslash{}{23\textbackslash{}}\textbackslash{}$. Formally: \textbackslash{}$\textbackslash{}$\textbackslash{}\textbackslash{}exists f \textbackslash{}\textbackslash{}in \textbackslash{}\textbackslash{}mathbb\textbackslash{}{Q\textbackslash{}}[X] \textbackslash{}\textbackslash{}text\textbackslash{}{ s.t. \textbackslash{}} \textbackslash{}\textbackslash{}text\textbackslash{}{Irreducible\textbackslash{}}(f) \textbackslash{}\textbackslash{}land \textbackslash{}\textbackslash{}deg(f) = 23 \textbackslash{}\textbackslash{}land \textbackslash{}\textbackslash{}text\textbackslash{}{Gal\textbackslash{}}(\textbackslash{}\textbackslash{}text\textbackslash{}{SplittingField\textbackslash{}}\textbackslash{}_\textbackslash{}{\textbackslash{}\textbackslash{}mathbb\textbackslash{}{Q\textbackslash{}}\textbackslash{}}(f) / \textbackslash{}\textbackslash{}mathbb\textbackslash{}{Q
\end{conjecture}

\begin{lstlisting}[caption={Lean 4 Sketch (v14 Galois)}]
import Mathlib.FieldTheory.IsGalois
import Mathlib.FieldTheory.SplittingField
import Mathlib.GroupTheory.SpecificGroups.Mathieu

open Polynomial

/--
Conjecture (Inverse Galois Problem for M₂₃): The Mathieu group M₂₃ can be realized as the
Galois group of a degree 23 irreducible polynomial over the rational numbers.
-/
theorem M23_is_Galois_over_Q :
  ∃ (f : ℚ[X]), Irreducible f ∧ f.natDegree = 23 ∧
    Nonempty (Gal (SplittingField f) ℚ ≃* (MathieuGroup 23)) := by
  sorry
  -- This is a major open problem. A proof is expected to proceed via the rigidity method.
  -- Step 1: Prove the existenc
\end{lstlisting}

\hrule\medskip

\subsection*{\texttt{hensley\textbackslash{}_hausdorff\textbackslash{}_dim}}

\textbf{Domain:} number\textbackslash{}_theory \quad \textbf{Class:} 4 \quad \textbf{Verdict:} \textcolor{incompleteamber}{\textbf{INCOMPLETE}}

\textbf{Confidence:} 0.70 \quad \textbf{Domain Prior:} 0.05 \quad \textbf{Sorry gaps:} 2

\textcolor{sorryred}{\textbf{Zero-Sorry Guillotine Applied:} 2 \texttt{sorry} stubs detected. Verdict downgraded from VERIFIED to INCOMPLETE.}

\begin{conjecture}
Let \textbackslash{}$d(A) := \textbackslash{}\textbackslash{}dim\textbackslash{}_H(\textbackslash{}\textbackslash{}\textbackslash{}{x \textbackslash{}\textbackslash{}in [0,1] \textbackslash{}\textbackslash{}setminus \textbackslash{}\textbackslash{}mathbb\textbackslash{}{Q\textbackslash{}} \textbackslash{}\textbackslash{}mid a\textbackslash{}_i(x) \textbackslash{}\textbackslash{}in A, \textbackslash{}\textbackslash{}forall i \textbackslash{}\textbackslash{}ge 1\textbackslash{}\textbackslash{}\textbackslash{}})\textbackslash{}$, where \textbackslash{}$a\textbackslash{}_i(x)\textbackslash{}$ are the partial quotients of the continued fraction of \textbackslash{}$x\textbackslash{}$ and \textbackslash{}$\textbackslash{}\textbackslash{}dim\textbackslash{}_H\textbackslash{}$ is the Hausdorff dimension. For \textbackslash{}$N \textbackslash{}\textbackslash{}in \textbackslash{}\textbackslash{}mathbb\textbackslash{}{N\textbackslash{}}\textbackslash{}_\textbackslash{}{\textbackslash{}\textbackslash{}ge 1\textbackslash{}}\textbackslash{}$, let \textbackslash{}$A\textbackslash{}_N = \textbackslash{}\textbackslash{}\textbackslash{}{1, 2, \textbackslash{}\textbackslash{}dots, N\textbackslash{}\textbackslash{}\textbackslash{}}\textbackslash{}$ and \textbackslash{}$B\textbackslash{}_N = \textbackslash{}\textbackslash{}\textbackslash{}{N+1, N+2, \textbackslash{}\textbackslash{}dots\textbackslash{}\textbackslash{}\textbackslash{}}\textbackslash{}$. Define the function \textbackslash{}$S(N) := d(A\textbackslash{}_N) + d(B\textbackslash{}_N)\textbackslash{}$. Then \textbackslash{}$S(N)\textbackslash{}$ is a strictly m
\end{conjecture}

\begin{lstlisting}[caption={Lean 4 Sketch (v14 Galois)}]
import Mathlib.NumberTheory.ContinuedFractions.Basic
import Mathlib.Analysis.Calculus.MeanValue
import Mathlib.MeasureTheory.Measure.Hausdorff

open Real Set Metric MeasureTheory

/-- The set of irrationals in [0,1] whose continued fraction digits are in a given set A. -/
noncomputable def ContinuedFraction.setOfDigits (A : Set ℕ) : Set ℝ :=
  {x | 0 ≤ x ∧ x ≤ 1 ∧ ¬ Irrational x ∧ ∀ i, (continuedFraction x).digits i ∈ A}

/-- The Hausdorff dimension of the set of numbers with continued fraction digits in A. -/
noncomputable def dimH_cf (A : Set ℕ) : ℝ := hausdorffDimension (ContinuedFraction.s
\end{lstlisting}

\hrule\medskip

\subsection*{\texttt{spherical\textbackslash{}_mode\textbackslash{}_quality\textbackslash{}_factor\textbackslash{}_te\textbackslash{}_tm}}

\textbf{Domain:} spectral\textbackslash{}_theory \quad \textbf{Class:} 3 \quad \textbf{Verdict:} \textcolor{incompleteamber}{\textbf{INCOMPLETE}}

\textbf{Confidence:} 0.70 \quad \textbf{Domain Prior:} 0.20 \quad \textbf{Sorry gaps:} 6

\textcolor{sorryred}{\textbf{Zero-Sorry Guillotine Applied:} 6 \texttt{sorry} stubs detected. Verdict downgraded from VERIFIED to INCOMPLETE.}

\begin{conjecture}
Let \textbackslash{}$z\textbackslash{}_\textbackslash{}{l,n\textbackslash{}}\textbackslash{}$ be the \textbackslash{}$n\textbackslash{}$-th positive zero of the spherical Bessel function of the first kind \textbackslash{}$j\textbackslash{}_l(x)\textbackslash{}$, for integers \textbackslash{}$l \textbackslash{}\textbackslash{}ge 1, n \textbackslash{}\textbackslash{}ge 1\textbackslash{}$. Define the total normalized excess loss for a given angular momentum number \textbackslash{}$l\textbackslash{}$ as the convergent series \textbackslash{}$S\textbackslash{}_l := \textbackslash{}\textbackslash{}\textbackslash{}\textbackslash{}sum\textbackslash{}_\textbackslash{}{n=1\textbackslash{}}\textbackslash{}textasciicircum{}\textbackslash{}\textbackslash{}\textbackslash{}\textbackslash{}infty \textbackslash{}\textbackslash{}\textbackslash{}\textbackslash{}frac\textbackslash{}{l(l+1)\textbackslash{}}\textbackslash{}{z\textbackslash{}_\textbackslash{}{l,n\textbackslash{}}\textbackslash{}textasciicircum{}2 - l(l+1)\textbackslash{}}\textbackslash{}$. Then the sequence \textbackslash{}$(S\textbackslash{}_l)\textbackslash{}_\textbackslash{}{l \textbackslash{}\textbackslash{}ge 1\textbackslash{}}\textbackslash{}$ is strictly increasing and satisfies the sharp bo
\end{conjecture}

\begin{lstlisting}[caption={Lean 4 Sketch (v14 Galois)}]
import Mathlib.Analysis.SpecialFunctions.Bessel\
import Mathlib.Topology.Instances.Real\
\
open Real SpecialFunctions\
\
-- We postulate the existence and properties of the n-th zero of the l-th spherical Bessel function.\
-- A full formalization would require developing the theory of zeros of these functions.\
axiom sphericalBessel_zero_exists (l n : ℕ) (hl : l ≥ 1) (hn : n ≥ 1): ∃! z, (sphericalBesselJ l z = 0) ∧ z > 0 ∧\
  (∀ y, 0 < y ∧ y < z → sphericalBesselJ l y ≠ 0 → ∀ m < n, Classical.choose (sphericalBessel_zero_exists l m hl (by linarith)) < y)\
\
-- Define z_{l,n} using the axiom\
d
\end{lstlisting}

\hrule\medskip


\newpage
%% ============================================================
\section{Open Problems and Future Work}
%% ============================================================

The following problems remain unresolved after the v14 re-run and represent
the frontier of formal AI mathematics:

\begin{itemize}
  \item \textbf{euler\_mascheroni\_closed\_form}: The irrationality of the
    Euler--Mascheroni constant $\gamma = \lim_{n\to\infty}(H_n - \ln n)$ is a
    major open problem. No closed-form algebraic expression is known.
  \item \textbf{feigenbaum\_alpha, feigenbaum\_delta}: The universality
    constants $\alpha \approx 2.502$ and $\delta \approx 4.669$ arise in
    period-doubling bifurcations. No algebraic proof of their transcendence exists.
  \item \textbf{anderson\_lyapunov\_exponent}: The Lyapunov exponent for
    the Anderson localization model requires deep results in random matrix theory.
\end{itemize}

\begin{thebibliography}{9}
\bibitem{horizonmath2024} Wang, E. et al.\ (2024). ``HorizonMath: A Benchmark
  for Open Mathematical Problems.'' arXiv preprint.
\bibitem{lean4} de Moura, L., Ullrich, S.\ (2021). ``The Lean 4 Theorem Prover
  and Programming Language.'' CADE-28.
\bibitem{mathlib4} The Mathlib Community (2024). ``Mathlib4.'' GitHub.
\bibitem{lakatos} Lakatos, I.\ (1976). \textit{Proofs and Refutations}.
  Cambridge University Press.
\end{thebibliography}

\end{document}
